% Options for packages loaded elsewhere
\PassOptionsToPackage{unicode}{hyperref}
\PassOptionsToPackage{hyphens}{url}
\documentclass[
  spanish,
]{article}
\usepackage{xcolor}
\usepackage{amsmath,amssymb}
\setcounter{secnumdepth}{-\maxdimen} % remove section numbering
\usepackage{iftex}
\ifPDFTeX
  \usepackage[T1]{fontenc}
  \usepackage[utf8]{inputenc}
  \usepackage{textcomp} % provide euro and other symbols
\else % if luatex or xetex
  \usepackage{unicode-math} % this also loads fontspec
  \defaultfontfeatures{Scale=MatchLowercase}
  \defaultfontfeatures[\rmfamily]{Ligatures=TeX,Scale=1}
\fi
\usepackage{lmodern}
\ifPDFTeX\else
  % xetex/luatex font selection
\fi
% Use upquote if available, for straight quotes in verbatim environments
\IfFileExists{upquote.sty}{\usepackage{upquote}}{}
\IfFileExists{microtype.sty}{% use microtype if available
  \usepackage[]{microtype}
  \UseMicrotypeSet[protrusion]{basicmath} % disable protrusion for tt fonts
}{}
\makeatletter
\@ifundefined{KOMAClassName}{% if non-KOMA class
  \IfFileExists{parskip.sty}{%
    \usepackage{parskip}
  }{% else
    \setlength{\parindent}{0pt}
    \setlength{\parskip}{6pt plus 2pt minus 1pt}}
}{% if KOMA class
  \KOMAoptions{parskip=half}}
\makeatother
\usepackage{longtable,booktabs,array}
\newcounter{none} % for unnumbered tables
\usepackage{calc} % for calculating minipage widths
% Correct order of tables after \paragraph or \subparagraph
\usepackage{etoolbox}
\makeatletter
\patchcmd\longtable{\par}{\if@noskipsec\mbox{}\fi\par}{}{}
\makeatother
% Allow footnotes in longtable head/foot
\IfFileExists{footnotehyper.sty}{\usepackage{footnotehyper}}{\usepackage{footnote}}
\makesavenoteenv{longtable}
\ifLuaTeX
\usepackage[bidi=basic,shorthands=off]{babel}
\else
\usepackage[bidi=default,shorthands=off]{babel}
\fi
\ifLuaTeX
  \usepackage{selnolig} % disable illegal ligatures
\fi
\setlength{\emergencystretch}{3em} % prevent overfull lines
\providecommand{\tightlist}{%
  \setlength{\itemsep}{0pt}\setlength{\parskip}{0pt}}
\DeclareUnicodeCharacter{22C5}{\ensuremath{\cdot}}
% PDF reproducible: sin /ID aleatorio (la fecha la fija SOURCE_DATE_EPOCH, ver scripts/comun.ps1).
\ifdefined\pdftrailerid\pdftrailerid{}\fi
\usepackage{bookmark}
\IfFileExists{xurl.sty}{\usepackage{xurl}}{} % add URL line breaks if available
\urlstyle{same}
\hypersetup{
  pdfauthor={Julián Calderón Almendros},
  pdflang={es},
  hidelinks,
  pdfcreator={LaTeX via pandoc}}

\ifXeTeX
\special{pdf:trailerid [ <b757066b316d6340143999b6560edc93> <b757066b316d6340143999b6560edc93> ]}
\fi
\ifPDFTeX
\pdftrailerid{}
\pdftrailer{/ID [ <b757066b316d6340143999b6560edc93> <b757066b316d6340143999b6560edc93> ]}
\fi
\ifLuaTeX
\pdfvariable trailerid {[ <b757066b316d6340143999b6560edc93> <b757066b316d6340143999b6560edc93> ]}
\fi

\author{Julián Calderón Almendros}
\date{}

\begin{document}

\section{Tema 1. Mecánica teórica}\label{tema-1.-mecuxe1nica-teuxf3rica}

\begin{itemize}
\tightlist
\item
  \textbf{Asignatura:} Complementos de Mecánica Cuántica (Máster en
  Computación Cuántica)
\item
  \textbf{Propósito:} apuntes de autoestudio elaborados para preparar la
  asignatura.
\item
  \textbf{Última edición:} 2026-09-24
\item
  \textbf{Autor:} Julián Calderón Almendros
\item
  \textbf{Correo electrónico:} julian.calderon.almendros at gmail.com
\item
  \textbf{GitHub:} \href{https://github.com/julian1c2a}{@julian1c2a}
\item
  \textbf{Proyecto:} \url{https://github.com/julian1c2a/MCC}
\item
  \textbf{Licencia:}
  \href{https://creativecommons.org/licenses/by/4.0/deed.es}{Creative
  Commons Reconocimiento 4.0 Internacional (CC BY 4.0)}. Se permite
  copiar, redistribuir, modificar y reutilizar este material con
  cualquier finalidad, incluso comercial, siempre que se reconozca la
  autoría original, se enlace la licencia y se indique si se han hecho
  cambios.
\end{itemize}

\subsection{1. Mecánica ondulatoria}\label{mecuxe1nica-ondulatoria}

La mecánica ondulatoria estudia cómo se propagan las perturbaciones a
través de un medio (o el vacío) en forma de ondas. Para comprender
matemáticamente esta propagación espacial y temporal, es útil comenzar
analizando el comportamiento de un único oscilador y luego expandirlo a
un sistema continuo.

\subsubsection{1.1. El oscilador simple: masa atada a un elástico (o
muelle) en el
techo}\label{el-oscilador-simple-masa-atada-a-un-eluxe1stico-o-muelle-en-el-techo}

Imaginemos un sistema básico: una masa \(m\) suspendida del techo
mediante un elástico o resorte ideal con una constante elástica \(k\).
Si tiramos de la masa hacia abajo separándola de su punto de equilibrio
y la soltamos, comenzará a oscilar verticalmente.

Este sistema describe un Movimiento Armónico Simple (MAS). Aplicando la
Segunda Ley de Newton y la Ley de Hooke, obtenemos la siguiente ecuación
diferencial para el movimiento:

\[
\begin{aligned}
m \cdot \dfrac{d^2y}{dt^2} = -k \cdot y
\end{aligned}
\]

La solución a esta ecuación nos da la posición \(y\) de la masa en
función de una única variable, el tiempo (\(t\)):

\[
\begin{aligned}
y(t) = A \cdot \cos(\omega \cdot t + \phi)
\end{aligned}
\]

En esta función:

\begin{itemize}
\tightlist
\item
  \(A\) es la amplitud (el desplazamiento máximo desde el equilibrio).
\item
  \(\omega\) es la frecuencia angular, relacionada con las propiedades
  del sistema (\(\omega = \sqrt{\dfrac{k}{m}}\)).
\item
  \(\phi\) es la constante de fase inicial.
\end{itemize}

Este ejemplo ilustra una oscilación local; la energía se transforma de
potencial a cinética, pero no se propaga a través del espacio.

\subsubsection{1.2. Movimiento ondulatorio: cuerda elástica horizontal
entre dos puntos
fijos}\label{movimiento-ondulatorio-cuerda-eluxe1stica-horizontal-entre-dos-puntos-fijos}

Para pasar al movimiento ondulatorio, consideremos ahora una cuerda
elástica tensada horizontalmente entre dos paredes (dos puntos fijos).
Podemos imaginar esta cuerda como una cadena infinita de pequeñas masas
conectadas por pequeños elásticos.

Si perturbamos un punto de la cuerda, esa oscilación no se queda en un
solo lugar (como la masa en el techo), sino que tira de los fragmentos
vecinos de la cuerda, haciendo que la perturbación viaje a lo largo del
eje \(x\).

En este caso, el desplazamiento transversal (la altura de la cuerda),
que llamaremos \(y\), ya no depende solo del momento en que lo miremos,
sino también de en qué parte de la cuerda nos fijemos. Por tanto, la
función es bivariable: \(y(x,t)\).

Una perturbación viajera (onda armónica) moviéndose a lo largo de la
cuerda se describe matemáticamente como:

\[
\begin{aligned}
y(x,t) = A \cdot \sin(k \cdot x \pm \omega \cdot t)
\end{aligned}
\]

Donde se introduce una nueva variable espacial:

\begin{itemize}
\tightlist
\item
  \(k\) es el número de onda (\(k = \dfrac{2 \cdot \pi}{\lambda}\)), que
  describe la periodicidad en el espacio (\(\lambda\) es la longitud de
  onda).
\end{itemize}

El signo del argumento determina el sentido de propagación. Con
\(k \cdot x - \omega \cdot t\), un punto de fase constante cumple
\(k \cdot x - \omega \cdot t = \text{cte}\), es decir,
\(x = \dfrac{\omega}{k} \cdot t + \text{cte}\): la perturbación avanza
hacia \(x\) crecientes. Con \(k \cdot x + \omega \cdot t\), el mismo
razonamiento da \(x = -\dfrac{\omega}{k} \cdot t + \text{cte}\): la
perturbación avanza hacia \(x\) decrecientes.

Nota sobre la notación: en el apartado 1.1 la letra \(k\) designaba la
constante elástica del muelle; a partir de aquí designa el número de
onda. Son magnitudes distintas (la primera se mide en N/m y la segunda
en rad/m), y el contexto indica en cada caso a cuál se refiere.

\paragraph{1.2.1. Ondas estacionarias}\label{ondas-estacionarias}

Si la cuerda está sujeta en dos puntos fijos, por ejemplo en \(x=0\) y
\(x=L\), una onda que viaja hacia un extremo se refleja en él y vuelve
en sentido contrario. En la cuerda coexisten entonces dos ondas de igual
amplitud y frecuencia que se propagan en sentidos opuestos:

\[
\begin{aligned}
y(x,t) = A \cdot \sin(k \cdot x - \omega \cdot t) + A \cdot \sin(k \cdot x + \omega \cdot t)
\end{aligned}
\]

Aplicando la identidad trigonométrica
\(\sin a + \sin b = 2 \cdot \sin\left(\dfrac{a+b}{2}\right) \cdot \cos\left(\dfrac{a-b}{2}\right)\)
con \(a = k \cdot x - \omega \cdot t\) y
\(b = k \cdot x + \omega \cdot t\), se obtiene
\(\dfrac{a+b}{2} = k \cdot x\) y \(\dfrac{a-b}{2} = -\omega \cdot t\).
Como el coseno es par, \(\cos(-\omega \cdot t) = \cos(\omega \cdot t)\),
y la suma queda:

\[
\begin{aligned}
y(x,t) = 2 \cdot A \cdot \sin(k \cdot x) \cdot \cos(\omega \cdot t)
\end{aligned}
\]

Esta superposición se llama onda estacionaria. La expresión es el
producto de un factor que solo depende de \(x\) y otro que solo depende
de \(t\). Su lectura es la siguiente: cada punto \(x\) de la cuerda
realiza un Movimiento Armónico Simple en el tiempo
(\(\cos(\omega \cdot t)\)), con una amplitud que depende de su posición
(\(2 \cdot A \cdot \sin(k \cdot x)\)).

Las condiciones de contorno (los extremos no pueden moverse) restringen
los valores posibles de \(k\):

\begin{itemize}
\tightlist
\item
  En \(x = 0\):
  \(y(0,t) = 2 \cdot A \cdot \sin(0) \cdot \cos(\omega \cdot t) = 0\)
  para todo \(t\). Esta condición se cumple automáticamente.
\item
  En \(x = L\): \(y(L,t) = 0\) para todo \(t\) exige
  \(\sin(k \cdot L) = 0\), es decir, \(k \cdot L = n \cdot \pi\) con
  \(n\) entero positivo.
\end{itemize}

Por tanto, solo son posibles los números de onda y longitudes de onda

\[
\begin{aligned}
k_n = \dfrac{n \cdot \pi}{L}, \qquad \lambda_n = \dfrac{2 \cdot \pi}{k_n} = \dfrac{2 \cdot L}{n}, \qquad n = 1, 2, 3, \dots
\end{aligned}
\]

Cada valor de \(n\) define un modo normal (o armónico) de la cuerda. El
valor \(n=0\) se excluye porque da \(y \equiv 0\).

Los puntos donde \(\sin(k \cdot x) = 0\) no se mueven nunca; se llaman
\textbf{nodos}. Para el modo \(n\) están en
\(x = \dfrac{j \cdot L}{n}\), con \(j = 0, 1, \dots, n\). Los puntos
donde \(|\sin(k \cdot x)| = 1\) oscilan con la amplitud máxima
\(2 \cdot A\); se llaman \textbf{vientres} o antinodos.

\subsubsection{1.3. Generalización: representación
compleja}\label{generalizaciuxf3n-representaciuxf3n-compleja}

En física es habitual escribir la función de onda en notación compleja.
Apoyándonos en la fórmula de Euler
(\(e^{i \cdot \theta} = \cos(\theta) + i \cdot \sin(\theta)\)), una onda
viajera puede expresarse en su forma compleja como:

\[
\begin{aligned}
y(x,t) = A \cdot e^{i \cdot (\omega \cdot t - k \cdot x)}
\end{aligned}
\]

Dado que las magnitudes físicas observables en mecánica clásica (como el
desplazamiento real de una cuerda) deben ser números reales, la onda
física se entiende simplemente como la parte real (o imaginaria) de esta
expresión compleja:

\[
\begin{aligned}
y_{real}(x,t) = \Re\{A \cdot e^{i \cdot (\omega \cdot t - k \cdot x)}\} = A \cdot \cos(\omega \cdot t - k \cdot x)
\end{aligned}
\]

Motivos para usar la notación compleja:

\begin{enumerate}
\def\labelenumi{\arabic{enumi}.}
\tightlist
\item
  Simplicidad del cálculo: al derivar o integrar una exponencial, la
  función conserva su forma y solo aparece un factor constante. Derivar
  respecto al tiempo equivale a multiplicar por \(i \cdot \omega\), y
  derivar respecto a la posición equivale a multiplicar por
  \(-i \cdot k\): \[
  \begin{aligned}
  \dfrac{\partial}{\partial t} e^{i \cdot (\omega \cdot t - k \cdot x)} = i \cdot \omega \cdot e^{i \cdot (\omega \cdot t - k \cdot x)}, \qquad
  \dfrac{\partial}{\partial x} e^{i \cdot (\omega \cdot t - k \cdot x)} = -i \cdot k \cdot e^{i \cdot (\omega \cdot t - k \cdot x)}
  \end{aligned}
  \] Con senos y cosenos, en cambio, cada derivada intercambia una
  función por la otra y cambia signos. Además, el producto de
  exponenciales se reduce a sumar exponentes
  (\(e^{a} \cdot e^{b} = e^{a+b}\)), lo que evita las identidades
  trigonométricas al superponer ondas.
\item
  Fase inicial: un desfase inicial \(\phi\) se incorpora definiendo una
  amplitud compleja \(\tilde{A} = A \cdot e^{i \cdot \phi}\), ya que
  \(\tilde{A} \cdot e^{i \cdot (\omega \cdot t - k \cdot x)} = A \cdot e^{i \cdot (\omega \cdot t - k \cdot x + \phi)}\).
\item
  Relación con la mecánica cuántica: en mecánica clásica la parte
  imaginaria es solo una herramienta de cálculo y al final se toma la
  parte real. En mecánica cuántica, en cambio, la función de onda de una
  partícula \(\Psi(x,t)\) es compleja por naturaleza y no se toma su
  parte real.
\end{enumerate}

Limitación del método: tomar la parte real conmuta con las operaciones
lineales (sumar, multiplicar por constantes reales, derivar, integrar),
pero no con el producto. En general,
\(\Re\{z_1 \cdot z_2\} \neq \Re\{z_1\} \cdot \Re\{z_2\}\). Por eso,
cuando se calculan magnitudes cuadráticas, como la energía o la
intensidad, hay que tomar primero la parte real de cada factor.

\subsubsection{1.4. La ecuación de ondas}\label{la-ecuaciuxf3n-de-ondas}

Cualquier perturbación física que se propague en forma de onda (ya sea
el desplazamiento de nuestra cuerda elástica, la presión del aire en el
sonido, o los campos eléctricos y magnéticos en la luz) debe obedecer a
una ecuación diferencial específica: la ecuación de ondas.

\paragraph{1.4.1. La ecuación diferencial en una
dimensión}\label{la-ecuaciuxf3n-diferencial-en-una-dimensiuxf3n}

Para una perturbación \(y(x,t)\) que viaja a lo largo de una sola
dimensión (el eje \(x\)), la ecuación clásica de ondas es una ecuación
diferencial en derivadas parciales, lineal y de segundo orden, expresada
como:

\[
\begin{aligned}
\dfrac{\partial^2 y}{\partial x^2} = \dfrac{1}{v^2} \cdot \dfrac{\partial^2 y}{\partial t^2}
\end{aligned}
\]

Aquí, \(v\) representa la velocidad de propagación (o velocidad de fase)
de la onda en el medio. Esta ecuación nos dice, fundamentalmente, que la
curvatura espacial de la onda en un punto (segunda derivada respecto a
\(x\)) es proporcional a la aceleración de ese punto (segunda derivada
respecto a \(t\)).

\paragraph{1.4.2. Comprobación con la representación
compleja}\label{comprobaciuxf3n-con-la-representaciuxf3n-compleja}

Se puede comprobar que la expresión compleja del apartado anterior,
\(y(x,t) = A \cdot e^{i \cdot (\omega \cdot t - k \cdot x)}\), es
solución de la ecuación de ondas. Como la derivada de la exponencial es
proporcional a la propia exponencial, las segundas derivadas parciales
se obtienen directamente:

\begin{enumerate}
\def\labelenumi{\arabic{enumi}.}
\tightlist
\item
  Derivada temporal: Al derivar dos veces respecto al tiempo,
  multiplicamos por \(i \cdot \omega\) cada vez. \[
  \begin{aligned}
  \dfrac{\partial^2 y}{\partial t^2} = (i \cdot \omega)^2 \cdot A \cdot e^{i \cdot (\omega \cdot t - k \cdot x)} = -\omega^2 \cdot y
  \end{aligned}
  \]
\item
  Derivada espacial: Al derivar dos veces respecto a la posición,
  multiplicamos por \(-i \cdot k\) cada vez. \[
  \begin{aligned}
  \dfrac{\partial^2 y}{\partial x^2} = (-i \cdot k)^2 \cdot A \cdot e^{i \cdot (\omega \cdot t - k \cdot x)} = -k^2 \cdot y
  \end{aligned}
  \]
\end{enumerate}

Si sustituimos estos resultados de nuevo en la ecuación de ondas
general, obtenemos:

\[
\begin{aligned}
-k^2 \cdot y = \dfrac{1}{v^2} \cdot (-\omega^2 \cdot y)
\end{aligned}
\]

Dividiendo entre \(-y\) (en los puntos donde \(y \neq 0\); como la
igualdad debe cumplirse para todo \(x\) y \(t\), basta con uno de
ellos), se obtiene una relación entre los parámetros de la onda:

\[
\begin{aligned}
k^2 = \dfrac{\omega^2}{v^2} \implies v = \dfrac{\omega}{k}
\end{aligned}
\]

donde se han tomado \(\omega\), \(k\) y \(v\) positivos. Por tanto, la
exponencial compleja es solución si y solo si \(\omega\) y \(k\)
satisfacen esta relación.

Antes de interpretar este resultado, conviene precisar dos conceptos.
Una \textbf{relación de dispersión} es una ecuación que vincula la
frecuencia angular \(\omega\) de una onda con su número de onda \(k\),
es decir, que indica cómo depende \(\omega\) de \(k\). En este caso, la
ecuación anterior puede escribirse como \(\omega = v \cdot k\): la
relación es lineal y la velocidad de fase \(v = \omega/k\) no depende de
la frecuencia. Por ello, las distintas componentes de frecuencia de un
pulso viajan todas a la misma velocidad y el pulso conserva su forma; a
estas ondas se las llama \textbf{no dispersivas}. Si la velocidad
dependiera de \(\omega\) o de \(k\), cada componente avanzaría a una
velocidad distinta y el pulso se deformaría al propagarse: habría
dispersión.

En resumen, \(\omega = v \cdot k\) es la relación de dispersión de las
ondas no dispersivas: la velocidad de la onda, \(v = \omega/k\), es la
misma para cualquier frecuencia angular \(\omega\) y cualquier número de
onda \(k\).

\paragraph{1.4.3. La solución de d'Alembert (ondas viajeras
arbitrarias)}\label{la-soluciuxf3n-de-dalembert-ondas-viajeras-arbitrarias}

Las soluciones armónicas (senos, cosenos o exponenciales complejas) son
un caso particular. Jean le Rond d'Alembert demostró que la ecuación de
ondas unidimensional admite soluciones de forma arbitraria.

Definamos los argumentos de las dos ondas viajeras como
\(u_- = \omega \cdot t - k \cdot x\) y
\(u_+ = \omega \cdot t + k \cdot x\). Si \(f\) y \(g\) son funciones
arbitrarias dos veces diferenciables, la solución general de d'Alembert
es:

\[
\begin{aligned}
y(x,t) = f(u_-) + g(u_+) = f(\omega \cdot t - k \cdot x) + g(\omega \cdot t + k \cdot x)
\end{aligned}
\]

Para comprobarlo se aplica la regla de la cadena. Las derivadas de los
argumentos son \(\dfrac{\partial u_\mp}{\partial t} = \omega\),
\(\dfrac{\partial u_-}{\partial x} = -k\) y
\(\dfrac{\partial u_+}{\partial x} = k\). Por tanto, las primeras
derivadas son:

\[
\begin{aligned}
\dfrac{\partial y}{\partial t} &= \omega \cdot f'(u_-) + \omega \cdot g'(u_+), \\
\dfrac{\partial y}{\partial x} &= -k \cdot f'(u_-) + k \cdot g'(u_+).
\end{aligned}
\]

Derivando una segunda vez, cada factor \(\pm k\) o \(\omega\) aparece al
cuadrado y los signos desaparecen:

\[
\begin{aligned}
\dfrac{\partial^2 y}{\partial t^2} &= \omega^2 \cdot \left[f''(u_-) + g''(u_+)\right], \\
\dfrac{\partial^2 y}{\partial x^2} &= k^2 \cdot \left[f''(u_-) + g''(u_+)\right].
\end{aligned}
\]

Como la relación de dispersión obtenida antes es \(\omega = v \cdot k\),
se cumple:

\[
\begin{aligned}
\dfrac{\partial^2 y}{\partial x^2}
&= \dfrac{\omega^2}{v^2} \cdot \left[f''(u_-) + g''(u_+)\right] \\
&= \dfrac{1}{v^2} \cdot \dfrac{\partial^2 y}{\partial t^2}.
\end{aligned}
\]

Por tanto, cualquier par de funciones \(f\) y \(g\) dos veces
diferenciables produce una solución. La primera componente,
\(f(\omega \cdot t - k \cdot x)\), viaja hacia la derecha: su argumento
es constante cuando
\(x = \dfrac{\omega}{k} \cdot t + \text{cte} = v \cdot t + \text{cte}\),
de modo que el perfil de \(f\) se traslada sin deformarse con velocidad
\(v\) hacia \(x\) crecientes. Por el mismo razonamiento,
\(g(\omega \cdot t + k \cdot x)\) viaja hacia la izquierda con velocidad
\(v\). Esto permite describir, por ejemplo, un pulso que se propaga y
otro que regresa tras reflejarse en un extremo de la cuerda.

\paragraph{1.4.4. Solución general por separación de
variables}\label{soluciuxf3n-general-por-separaciuxf3n-de-variables}

La forma más sistemática de resolver la ecuación de ondas (especialmente
cuando hay condiciones de contorno fijas, como en nuestra cuerda) es el
método de separación de variables. Asumimos como hipótesis que la
solución \(y(x,t)\) puede expresarse como el producto de dos funciones
independientes, una que depende solo de la posición \(X(x)\) y otra solo
del tiempo \(T(t)\):

\[
\begin{aligned}
y(x,t) = X(x) \cdot T(t)
\end{aligned}
\]

Si calculamos las segundas derivadas y las sustituimos en la ecuación de
ondas original, obtenemos:

\[
\begin{aligned}
X''(x) \cdot T(t) = \dfrac{1}{v^2} \cdot X(x) \cdot T''(t)
\end{aligned}
\]

Dividiendo toda la ecuación por \(X(x) \cdot T(t)\), separamos las
variables a cada lado de la igualdad:

\[
\begin{aligned}
\dfrac{X''(x)}{X(x)} = \dfrac{1}{v^2} \cdot \dfrac{T''(t)}{T(t)}
\end{aligned}
\]

Dado que el lado izquierdo depende exclusivamente de \(x\) y el lado
derecho exclusivamente de \(t\), la única forma de que esta igualdad se
cumpla para cualquier valor de \(x\) y \(t\) es que ambos lados sean
iguales a una misma constante. (Si se fija \(t\) y se varía \(x\), el
lado derecho no cambia, luego el izquierdo tampoco puede cambiar; y
viceversa.)

El signo de la constante de separación lo deciden las condiciones de
contorno de la cuerda fija, \(X(0) = X(L) = 0\):

\begin{itemize}
\tightlist
\item
  Si la constante es positiva, \(+\mu^2\), la ecuación
  \(X'' = \mu^2 \cdot X\) tiene soluciones
  \(X = C_1 \cdot e^{\mu \cdot x} + C_2 \cdot e^{-\mu \cdot x}\).
  Imponer \(X(0)=0\) da \(C_2 = -C_1\), y entonces
  \(X(L) = 2 \cdot C_1 \cdot \sinh(\mu \cdot L) = 0\) obliga a
  \(C_1 = 0\): solo queda la solución trivial.
\item
  Si la constante es cero, \(X'' = 0\) da \(X = C_1 + C_2 \cdot x\), y
  las dos condiciones de contorno obligan a \(C_1 = C_2 = 0\).
\item
  Si la constante es negativa, las soluciones son senos y cosenos, que
  sí pueden anularse en ambos extremos.
\end{itemize}

Por tanto, la constante de separación debe ser negativa, y la llamaremos
\(-k^2\):

\[
\begin{aligned}
\dfrac{X''(x)}{X(x)} = -k^2 \quad \text{y} \quad \dfrac{1}{v^2} \cdot \dfrac{T''(t)}{T(t)} = -k^2
\end{aligned}
\]

Esto descompone la ecuación en derivadas parciales en dos Ecuaciones
Diferenciales Ordinarias (EDOs) simples:

\begin{enumerate}
\def\labelenumi{\arabic{enumi}.}
\tightlist
\item
  Ecuación espacial: \(X''(x) + k^2 \cdot X(x) = 0\)
\item
  Ecuación temporal: \(T''(t) + \omega^2 \cdot T(t) = 0 \quad\) (donde
  hemos definido \(\omega = k \cdot v\))
\end{enumerate}

La solución general para estas ecuaciones son combinaciones lineales de
senos y cosenos (o exponenciales complejas):

\[
\begin{aligned}
X(x) = A_1 \cdot \cos(k \cdot x) + B_1 \cdot \sin(k \cdot x)
\end{aligned}
\] \[
\begin{aligned}
T(t) = A_2 \cdot \cos(\omega \cdot t) + B_2 \cdot \sin(\omega \cdot t)
\end{aligned}
\]

Por lo tanto, la solución general para un modo normal de vibración es el
producto de ambas:

\[
\begin{aligned}
y(x,t) = \left( A_1 \cdot \cos(k \cdot x) + B_1 \cdot \sin(k \cdot x) \right) \cdot \left( A_2 \cdot \cos(\omega \cdot t) + B_2 \cdot \sin(\omega \cdot t) \right)
\end{aligned}
\]

Las condiciones de contorno fijan varias de estas constantes:

\begin{itemize}
\tightlist
\item
  \(y(0,t) = 0\) para todo \(t\) exige \(X(0) = A_1 = 0\), de modo que
  \(X(x) = B_1 \cdot \sin(k \cdot x)\).
\item
  \(y(L,t) = 0\) para todo \(t\) exige
  \(B_1 \cdot \sin(k \cdot L) = 0\). Como \(B_1 = 0\) daría la solución
  trivial, debe ser \(\sin(k \cdot L) = 0\), es decir,
  \(k_n = \dfrac{n \cdot \pi}{L}\) y, en consecuencia,
  \(\omega_n = v \cdot k_n = \dfrac{n \cdot \pi \cdot v}{L}\).
\end{itemize}

Absorbiendo \(B_1\) en las constantes temporales, el modo \(n\) es
\(y_n(x,t) = \sin(k_n \cdot x) \cdot \left( a_n \cdot \cos(\omega_n \cdot t) + b_n \cdot \sin(\omega_n \cdot t) \right)\),
que es la onda estacionaria del apartado 1.2.1. Como la ecuación de
ondas es lineal, la suma de todos los modos también es solución; esa
suma es una serie de Fourier en senos, y sus coeficientes \(a_n\) y
\(b_n\) se determinan a partir de la forma y la velocidad iniciales de
la cuerda, del mismo modo que se hace en el apartado siguiente para la
membrana.

\paragraph{1.4.5. Generalización a dos dimensiones: membranas y series
de
Fourier}\label{generalizaciuxf3n-a-dos-dimensiones-membranas-y-series-de-fourier}

Para una superficie, como una membrana tensa, el desplazamiento
transversal \(y(x,z,t)\) depende de dos coordenadas espaciales. La
ecuación de ondas bidimensional es:

\[
\begin{aligned}
\dfrac{\partial^2 y}{\partial t^2}
= v^2 \cdot \left( \dfrac{\partial^2 y}{\partial x^2} + \dfrac{\partial^2 y}{\partial z^2} \right).
\end{aligned}
\]

En este caso, la fórmula de d'Alembert no proporciona la solución
general: una perturbación puede propagarse en infinitas direcciones
sobre la superficie. Para una membrana rectangular con lados \(L_x\) y
\(L_z\) fijados, \(y(0,z,t)=y(L_x,z,t)=y(x,0,t)=y(x,L_z,t)=0\), sí se
puede construir la solución general mediante una \textbf{serie doble de
Fourier}:

\[
\begin{aligned}
y(x,z,t)
&= \sum_{m=1}^{\infty} \sum_{n=1}^{\infty}
\left[A_{mn} \cdot \cos(\omega_{mn} \cdot t) + B_{mn} \cdot \sin(\omega_{mn} \cdot t)\right] \\
&\quad \cdot \sin\left(\dfrac{m \cdot \pi \cdot x}{L_x}\right)
\cdot \sin\left(\dfrac{n \cdot \pi \cdot z}{L_z}\right).
\end{aligned}
\]

Cada par de enteros positivos \((m,n)\) identifica un modo normal de
vibración. Las frecuencias angulares permitidas son:

\[
\begin{aligned}
\omega_{mn} = v \cdot \pi \cdot \sqrt{\left(\dfrac{m}{L_x}\right)^2 + \left(\dfrac{n}{L_z}\right)^2}.
\end{aligned}
\]

Para determinar los coeficientes, fijamos el desplazamiento inicial
\(y_0(x,z)=y(x,z,0)\) y la velocidad inicial
\(u_0(x,z)=\left.\dfrac{\partial y}{\partial t}\right|_{t=0}\). La
ortogonalidad de las funciones seno en el rectángulo da directamente:

\[
\begin{aligned}
A_{mn}
= \dfrac{4}{L_x \cdot L_z}
\int_0^{L_x}\int_0^{L_z}
y_0(x,z)
\cdot \sin\left(\dfrac{m \cdot \pi \cdot x}{L_x}\right)
\cdot \sin\left(\dfrac{n \cdot \pi \cdot z}{L_z}\right)
\, dz\, dx,
\\
B_{mn}
= \dfrac{4}{L_x \cdot L_z \cdot \omega_{mn}}
\int_0^{L_x}\int_0^{L_z}
u_0(x,z)
\cdot \sin\left(\dfrac{m \cdot \pi \cdot x}{L_x}\right)
\cdot \sin\left(\dfrac{n \cdot \pi \cdot z}{L_z}\right)
\, dz\, dx.
\end{aligned}
\]

La primera fórmula es la proyección del desplazamiento inicial sobre el
modo \((m,n)\). Para la segunda, se deriva la expansión respecto del
tiempo, se evalúa en \(t=0\) y se proyecta la velocidad inicial; de ahí
el factor adicional \(1/\omega_{mn}\). Por tanto, las series de Fourier
no son un recurso adicional opcional en este problema: son el mecanismo
que combina todos los modos normales para reproducir una forma inicial
arbitraria compatible con los bordes.

\paragraph{1.4.6. Generalización a tres
dimensiones}\label{generalizaciuxf3n-a-tres-dimensiones}

Si la perturbación se propaga en el espacio tridimensional (como el
sonido en una habitación o la luz de una estrella), la función dependerá
de tres coordenadas espaciales y el tiempo: \(\Psi(x,y,z,t)\). En este
caso, la derivada espacial de la ecuación unidimensional se sustituye
por el operador Laplaciano (\(\nabla^2\) o \(\Delta\)):

\[
\begin{aligned}
\nabla^2 \Psi = \dfrac{1}{v^2} \cdot \dfrac{\partial^2 \Psi}{\partial t^2}
\end{aligned}
\]

Donde el operador Laplaciano en coordenadas cartesianas se define como
la suma de las segundas derivadas parciales espaciales:
\(\nabla^2 = \dfrac{\partial^2}{\partial x^2} + \dfrac{\partial^2}{\partial y^2} + \dfrac{\partial^2}{\partial z^2}\).

\subsection{2. Modos normales y expansiones de Fourier en tres
dimensiones}\label{modos-normales-y-expansiones-de-fourier-en-tres-dimensiones}

\subsubsection{2.1. Modos normales en una cavidad
rectangular}\label{modos-normales-en-una-cavidad-rectangular}

En una región tridimensional acotada, las condiciones de contorno
seleccionan un conjunto discreto de modos normales. Consideremos una
cavidad rectangular de lados \(L_x\), \(L_y\) y \(L_z\), con el campo
\(\Psi\) anulado en todas sus paredes. Cada modo espacial permitido
tiene la forma:

\[
\begin{aligned}
\Phi_{mnp}(x,y,z)
= \sin\left(\dfrac{m \cdot \pi \cdot x}{L_x}\right)
\cdot \sin\left(\dfrac{n \cdot \pi \cdot y}{L_y}\right)
\cdot \sin\left(\dfrac{p \cdot \pi \cdot z}{L_z}\right),
\end{aligned}
\]

donde \(m\), \(n\) y \(p\) son enteros positivos. Al combinar esta parte
espacial con una oscilación temporal, la frecuencia angular de cada modo
es:

\[
\begin{aligned}
\omega_{mnp}
= v \cdot \pi \cdot \sqrt{
\left(\dfrac{m}{L_x}\right)^2
+ \left(\dfrac{n}{L_y}\right)^2
+ \left(\dfrac{p}{L_z}\right)^2}.
\end{aligned}
\]

Cada terna \((m,n,p)\) representa un patrón de vibración independiente,
con nodos determinados por las paredes y por los ceros de las funciones
seno. Estos modos son los análogos tridimensionales de los armónicos de
una cuerda fija.

\subsubsection{2.2. Expansión triple de Fourier y condiciones
iniciales}\label{expansiuxf3n-triple-de-fourier-y-condiciones-iniciales}

La solución general se obtiene superponiendo todos los modos normales:

\[
\begin{aligned}
\Psi(x,y,z,t)
&= \sum_{m=1}^{\infty} \sum_{n=1}^{\infty} \sum_{p=1}^{\infty}
\left[A_{mnp} \cdot \cos(\omega_{mnp} \cdot t)
+ B_{mnp} \cdot \sin(\omega_{mnp} \cdot t)\right] \\
&\quad \cdot \Phi_{mnp}(x,y,z).
\end{aligned}
\]

Para expresarlos de forma explícita, definimos el campo inicial
\(\Psi_0(x,y,z)=\Psi(x,y,z,0)\) y la velocidad inicial
\(V_0(x,y,z)=\left.\dfrac{\partial\Psi}{\partial t}\right|_{t=0}\). La
ortogonalidad de los modos seno sobre la cavidad da:

\[
\begin{aligned}
A_{mnp} &= \dfrac{8}{L_x \cdot L_y \cdot L_z}
\int_0^{L_x}\int_0^{L_y}\int_0^{L_z}
\Psi_0(x,y,z) \\
&\quad \cdot \sin\left(\dfrac{m \cdot \pi \cdot x}{L_x}\right)
\cdot \sin\left(\dfrac{n \cdot \pi \cdot y}{L_y}\right) \\
&\quad \cdot
\sin\left(\dfrac{p \cdot \pi \cdot z}{L_z}\right)
\, dz\, dy\, dx,
\\
B_{mnp} &= \dfrac{8}{L_x \cdot L_y \cdot L_z \cdot \omega_{mnp}}
\int_0^{L_x}\int_0^{L_y}\int_0^{L_z}
V_0(x,y,z) \\
&\quad \cdot \sin\left(\dfrac{m \cdot \pi \cdot x}{L_x}\right)
\cdot \sin\left(\dfrac{n \cdot \pi \cdot y}{L_y}\right) \\
&\quad \cdot
\sin\left(\dfrac{p \cdot \pi \cdot z}{L_z}\right)
\, dz\, dy\, dx.
\end{aligned}
\]

La primera integral proyecta el campo inicial sobre el modo \((m,n,p)\).
La segunda se obtiene derivando la expansión temporal, evaluándola en
\(t=0\) y proyectando la velocidad inicial; por eso incorpora el factor
\(1/\omega_{mnp}\). Cada coeficiente queda así determinado de manera
independiente.

Por tanto, las expansiones de Fourier convierten el problema de una
ecuación diferencial parcial en una familia infinita de osciladores
armónicos independientes, uno por cada modo normal. En dominios con otra
geometría se usan las funciones propias del Laplaciano adecuadas:
funciones de Bessel en cavidades cilíndricas o esféricas, por ejemplo.

\subsubsection{2.3. Geometría y elección de
coordenadas}\label{geometruxeda-y-elecciuxf3n-de-coordenadas}

Las coordenadas naturales no se eligen por la dimensión del problema,
sino por la geometría del dominio y de sus condiciones de contorno. Para
una membrana rectangular o una cavidad paralelepipédica, las coordenadas
cartesianas y las series de senos son las más adecuadas. En cambio, una
frontera circular o cilíndrica sugiere coordenadas polares o
cilíndricas; una frontera esférica sugiere coordenadas esféricas. En
cada caso, las funciones propias del Laplaciano se adaptan a los bordes
y forman la base de la expansión modal.

\paragraph{2.3.1. Membrana circular: coordenadas polares y series de
Fourier-Bessel}\label{membrana-circular-coordenadas-polares-y-series-de-fourier-bessel}

Para una membrana circular, con coordenadas polares \((r,\theta)\), el
Laplaciano es:

\[
\begin{aligned}
\nabla^2 y
= \dfrac{\partial^2 y}{\partial r^2}
+ \dfrac{1}{r} \cdot \dfrac{\partial y}{\partial r}
+ \dfrac{1}{r^2} \cdot \dfrac{\partial^2 y}{\partial \theta^2}.
\end{aligned}
\]

La ecuación de ondas sigue siendo
\(\dfrac{\partial^2 y}{\partial t^2}=v^2 \cdot \nabla^2y\), pero la
separación de variables ya no genera senos en la coordenada radial, sino
funciones de Bessel. Si el borde circular de radio \(R\) está fijo, los
modos espaciales tienen la forma:

\[
\begin{aligned}
J_m\left(\dfrac{\alpha_{mn} \cdot r}{R}\right)
\cdot \left[C_{mn} \cdot \cos(m \cdot \theta) + D_{mn} \cdot \sin(m \cdot \theta)\right],
\end{aligned}
\]

donde \(J_m\) es la función de Bessel de primera especie y
\(\alpha_{mn}\) es el enésimo cero positivo de \(J_m\). La condición
\(y(R,\theta,t)=0\) exige precisamente \(J_m(\alpha_{mn})=0\). La
frecuencia de cada modo es:

\[
\begin{aligned}
\omega_{mn} = \dfrac{v \cdot \alpha_{mn}}{R}.
\end{aligned}
\]

Una superposición de estos modos, con coeficientes temporales de seno y
coseno, constituye una expansión de Fourier-Bessel. Es el análogo
circular de la serie doble de Fourier usada para una membrana
rectangular.

\paragraph{2.3.2. Cavidad esférica: armónicos esféricos y Bessel
esféricas}\label{cavidad-esfuxe9rica-armuxf3nicos-esfuxe9ricos-y-bessel-esfuxe9ricas}

Para una cavidad esférica se usan coordenadas \((r,\theta,\varphi)\). El
Laplaciano adopta la forma:

\[
\begin{aligned}
\nabla^2 \Psi
= \dfrac{1}{r^2} \cdot \dfrac{\partial}{\partial r}
\left(r^2 \cdot \dfrac{\partial \Psi}{\partial r}\right)
+ \dfrac{1}{r^2 \cdot \sin\theta} \cdot \dfrac{\partial}{\partial\theta}
\left(\sin\theta \cdot \dfrac{\partial\Psi}{\partial\theta}\right)
+ \dfrac{1}{r^2 \cdot \sin^2\theta} \cdot \dfrac{\partial^2\Psi}{\partial\varphi^2}.
\end{aligned}
\]

La parte angular de los modos separados está descrita por los armónicos
esféricos \(Y_\ell^m(\theta,\varphi)\), mientras que la parte radial
regular en el origen está descrita por las funciones de Bessel esféricas
\(j_\ell(k \cdot r)\). Para una esfera de radio \(R\) con frontera fija,
una base modal es:

\[
\begin{aligned}
j_\ell\left(\dfrac{\beta_{\ell n} \cdot r}{R}\right)
\cdot Y_\ell^m(\theta,\varphi),
\end{aligned}
\]

donde \(\beta_{\ell n}\) es el enésimo cero positivo de \(j_\ell\). Las
frecuencias permitidas son \(\omega_{\ell n}=v \cdot \beta_{\ell n}/R\).
La solución se obtiene sumando estos modos con amplitudes determinadas
por el campo y la velocidad iniciales.

\paragraph{2.3.3. Espacio ilimitado: ondas planas y transformada de
Fourier}\label{espacio-ilimitado-ondas-planas-y-transformada-de-fourier}

Cuando no hay fronteras que impongan modos discretos, las ondas planas
son la base más simple. Una componente plana en cualquier dimensión
tiene la forma:

\[
\begin{aligned}
\Psi(\mathbf{r},t)
= A \cdot e^{i \cdot (\mathbf{k} \cdot \mathbf{r}-\omega \cdot t)},
\qquad \omega = v \cdot |\mathbf{k}|.
\end{aligned}
\]

La solución general se expresa entonces mediante una transformada de
Fourier, es decir, una superposición continua de vectores de onda
\(\mathbf{k}\) en lugar de una suma discreta de modos. Así, las series
de Fourier son apropiadas para dominios acotados, mientras que la
transformada de Fourier es la formulación natural en espacio ilimitado.

\subsection{3. Funcionales y cálculo de
variaciones}\label{funcionales-y-cuxe1lculo-de-variaciones}

\subsubsection{3.1. Funcionales y optimización
variacional}\label{funcionales-y-optimizaciuxf3n-variacional}

El cálculo de variaciones estudia problemas en los que la incógnita no
es un número (el valor de una variable que optimiza una función), sino
una función completa. Para ello se usa un \textbf{funcional}: una regla
que recibe una función completa como entrada y devuelve un número real.
Por ejemplo, una expresión de la forma

\[
\begin{aligned}
J[y] = \int_a^b F(x, y(x), y'(x)) \, dx
\end{aligned}
\]

asigna un número \(J[y]\) a cada curva admisible \(y(x)\). Las
condiciones de contorno, como fijar \(y(a)\) e \(y(b)\), delimitan qué
funciones pueden considerarse.

El problema de optimización variacional consiste en encontrar la función
admisible para la cual el funcional alcanza un valor mínimo, máximo o,
más generalmente, estacionario. El último caso es el análogo funcional
de imponer que una derivada ordinaria sea cero.

\paragraph{3.1.1. El problema de optimización variacional: funcionales
frente a
funciones}\label{el-problema-de-optimizaciuxf3n-variacional-funcionales-frente-a-funciones}

Antes de aplicar estas ideas a la mecánica, conviene distinguir con
precisión entre funciones y funcionales. A menudo se dice coloquialmente
que ``minimizamos la función \(f\) con respecto a \(\mathcal{L}\)'',
pero esto es una imprecisión.

En el cálculo diferencial ordinario, optimizamos una función (ej.
\(y(x)\)): buscamos un número \(x\) que nos devuelva el valor mínimo de
\(y\), haciendo \(\dfrac{dy}{dx} = 0\).

En el cálculo de variaciones, sin embargo, optimizamos un funcional (en
este caso, la Acción \(S\)). Un funcional toma como argumento una
función entera y devuelve un número real. El espacio de búsqueda no es
la recta de los números reales, sino un espacio de dimensión infinita
que contiene todas las posibles formas que podría adoptar el campo
\(f(x,t)\) entre las fronteras fijadas.

No minimizamos la función \(f\) punto a punto; buscamos cuál es el
``camino'' o ``superficie'' completa que hace estacionario el número
global producido por el funcional. En mecánica, ese funcional recibirá
el nombre de acción y se aplicará el Principio de Hamilton:
\(\delta S = 0\).

\paragraph{3.1.2. Primera variación de un funcional y ecuación de
Euler-Lagrange}\label{primera-variaciuxf3n-de-un-funcional-y-ecuaciuxf3n-de-euler-lagrange}

Veamos de dónde surge la condición que debe satisfacer una curva
estacionaria para el funcional general:

\[
\begin{aligned}
J[y] = \int_a^b F(x,y(x),y'(x))\,dx.
\end{aligned}
\]

Partimos de una curva candidata \(y(x)\) cuyos extremos están fijados.
Para compararla con curvas próximas, elegimos una función suave
arbitraria \(\eta(x)\) que se anula en los extremos,
\(\eta(a)=\eta(b)=0\), e introducimos una familia de variaciones:

\[
\begin{aligned}
y_\epsilon(x) = y(x) + \epsilon \cdot \eta(x),
\qquad
y_\epsilon'(x) = y'(x) + \epsilon \cdot \eta'(x),
\end{aligned}
\]

donde \(\epsilon\) es un número real pequeño. Ahora el funcional deja de
depender directamente de una función y pasa a ser una función ordinaria
del parámetro \(\epsilon\):

\[
\begin{aligned}
\Phi(\epsilon)
= J[y_\epsilon]
= \int_a^b F\left(x,y(x)+\epsilon \cdot \eta(x),y'(x)+\epsilon \cdot \eta'(x)\right)\,dx.
\end{aligned}
\]

Si \(y\) hace estacionario el funcional, entonces \(\Phi\) debe ser
estacionaria en \(\epsilon=0\). La primera variación se define como
\(\delta J=\Phi'(0)\); por tanto, la condición de estacionariedad es
\(\delta J=0\). Derivando bajo el signo integral y aplicando la regla de
la cadena (el integrando depende de \(\epsilon\) a través de su segundo
argumento, con derivada \(\eta\), y de su tercer argumento, con derivada
\(\eta'\)) se obtiene:

\[
\begin{aligned}
\delta J
&= \left.\dfrac{d\Phi}{d\epsilon}\right|_{\epsilon=0} \\
&= \int_a^b
\left[
\dfrac{\partial F}{\partial y} \cdot \eta
+ \dfrac{\partial F}{\partial y'} \cdot \eta'
\right]\,dx.
\end{aligned}
\]

El segundo término contiene la derivada de la variación. Lo reescribimos
mediante integración por partes:

\[
\begin{aligned}
\int_a^b \dfrac{\partial F}{\partial y'} \cdot \eta'\,dx
= \left[\dfrac{\partial F}{\partial y'} \cdot \eta\right]_a^b
- \int_a^b \dfrac{d}{dx}\left(\dfrac{\partial F}{\partial y'}\right) \cdot \eta\,dx.
\end{aligned}
\]

El término de borde es nulo porque \(\eta(a)=\eta(b)=0\). Al sustituirlo
en la primera variación queda:

\[
\begin{aligned}
\delta J
= \int_a^b
\left[
\dfrac{\partial F}{\partial y}
- \dfrac{d}{dx}\left(\dfrac{\partial F}{\partial y'}\right)
\right] \cdot \eta(x)\,dx.
\end{aligned}
\]

La función \(\eta(x)\) puede escogerse arbitrariamente en el interior
del intervalo. Por el lema fundamental del cálculo de variaciones, la
única forma de que esta integral sea cero para toda variación admisible
es que el corchete se anule punto a punto:

\[
\begin{aligned}
\dfrac{\partial F}{\partial y}
- \dfrac{d}{dx}\left(\dfrac{\partial F}{\partial y'}\right) = 0.
\end{aligned}
\]

Esta es la ecuación de Euler-Lagrange unidimensional. Es la expresión
general de la condición \(\delta J=0\) y será la ecuación que
aplicaremos al funcional de tiempo de la braquistócrona.

\paragraph{\texorpdfstring{3.1.3. Extensión a campos: funcionales de
\(f(x,t)\)}{3.1.3. Extensión a campos: funcionales de f(x,t)}}\label{extensiuxf3n-a-campos-funcionales-de-fxt}

En los problemas de ondas no buscamos una curva \(y(x)\), sino un campo
\(f(x,t)\) que asigna un valor a cada punto del espacio-tiempo. El
análogo bidimensional del funcional anterior integra sobre una región
\(\Omega\) del plano \((x,t)\):

\[
\begin{aligned}
S[f]
= \int_{t_1}^{t_2}\int_{x_1}^{x_2}
\mathcal{L}\left(
f,
\dfrac{\partial f}{\partial t},
\dfrac{\partial f}{\partial x},
x,t
\right)\,dx\,dt.
\end{aligned}
\]

Este funcional se llama \textbf{acción} cuando \(\mathcal{L}\) es una
densidad lagrangiana. Para variar el campo, elegimos una perturbación
suave \(\eta(x,t)\) que se anula en el borde de la región \(\Omega\), y
definimos:

\[
\begin{aligned}
f_\epsilon(x,t) = f(x,t) + \epsilon \cdot \eta(x,t).
\end{aligned}
\]

Sus derivadas varían de forma coherente:

\[
\begin{aligned}
\dfrac{\partial f_\epsilon}{\partial t}
= \dfrac{\partial f}{\partial t}
+ \epsilon \cdot \dfrac{\partial\eta}{\partial t},
\\
\dfrac{\partial f_\epsilon}{\partial x}
= \dfrac{\partial f}{\partial x}
+ \epsilon \cdot \dfrac{\partial\eta}{\partial x}.
\end{aligned}
\]

Como antes, introducimos \(\Phi(\epsilon)=S[f_\epsilon]\) y exigimos que
\(\Phi'(0)=\delta S=0\). La regla de la cadena aplicada al integrando
proporciona:

\[
\begin{aligned}
\delta S
= \int_{t_1}^{t_2}\int_{x_1}^{x_2}
\Bigg[
\dfrac{\partial\mathcal{L}}{\partial f} \cdot \eta
+ \dfrac{\partial\mathcal{L}}{\partial\left(\dfrac{\partial f}{\partial t}\right)}
\cdot \dfrac{\partial\eta}{\partial t}
+ \dfrac{\partial\mathcal{L}}{\partial\left(\dfrac{\partial f}{\partial x}\right)}
\cdot \dfrac{\partial\eta}{\partial x}
\Bigg]\,dx\,dt.
\end{aligned}
\]

Integramos por partes los dos términos que contienen derivadas de
\(\eta\). Los términos de borde se anulan porque la variación es nula en
el contorno. Por tanto:

\[
\begin{aligned}
\delta S
= \int_{t_1}^{t_2}\int_{x_1}^{x_2}
\Bigg[
\dfrac{\partial\mathcal{L}}{\partial f}
- \dfrac{\partial}{\partial t}
\left(\dfrac{\partial\mathcal{L}}{\partial\left(\dfrac{\partial f}{\partial t}\right)}\right)
- \dfrac{\partial}{\partial x}
\left(\dfrac{\partial\mathcal{L}}{\partial\left(\dfrac{\partial f}{\partial x}\right)}\right)
\Bigg] \cdot \eta(x,t)\,dx\,dt.
\end{aligned}
\]

Puesto que \(\eta(x,t)\) es arbitraria en el interior de \(\Omega\), el
lema fundamental conduce a la ecuación de Euler-Lagrange para un campo
en una dimensión espacial:

\[
\begin{aligned}
\dfrac{\partial\mathcal{L}}{\partial f}
- \dfrac{\partial}{\partial t}
\left(\dfrac{\partial\mathcal{L}}{\partial\left(\dfrac{\partial f}{\partial t}\right)}\right)
- \dfrac{\partial}{\partial x}
\left(\dfrac{\partial\mathcal{L}}{\partial\left(\dfrac{\partial f}{\partial x}\right)}\right)
= 0.
\end{aligned}
\]

La expresión unidimensional anterior se recupera si el campo depende
sólo de \(x\). En más dimensiones espaciales aparece un término análogo
por cada derivada espacial del campo.

\paragraph{3.1.4. Ecuación de Euler-Lagrange en tres dimensiones
espaciales}\label{ecuaciuxf3n-de-euler-lagrange-en-tres-dimensiones-espaciales}

Si el campo depende de las tres coordenadas espaciales y del tiempo,
\(f=f(x,y,z,t)\), la acción se integra sobre un volumen espacial y un
intervalo temporal:

\[
\begin{aligned}
S[f]
= \int_{t_1}^{t_2}\int_V
\mathcal{L}\left(
f,
\dfrac{\partial f}{\partial t},
\dfrac{\partial f}{\partial x},
\dfrac{\partial f}{\partial y},
\dfrac{\partial f}{\partial z},
x,y,z,t
\right)\,dV\,dt.
\end{aligned}
\]

A la vista del resultado para \(x\) y \(t\), no hace falta repetir toda
la deducción: se integra por partes una vez por cada derivada de la
variación. Si \(\eta\) se anula en la frontera espacial y en los
instantes extremos, los cuatro términos de borde desaparecen. La
ecuación resultante es:

\[
\begin{aligned}
\dfrac{\partial\mathcal{L}}{\partial f}
&- \dfrac{\partial}{\partial t}
\left(\dfrac{\partial\mathcal{L}}{\partial\left(\dfrac{\partial f}{\partial t}\right)}\right)
- \dfrac{\partial}{\partial x}
\left(\dfrac{\partial\mathcal{L}}{\partial\left(\dfrac{\partial f}{\partial x}\right)}\right) \\
&- \dfrac{\partial}{\partial y}
\left(\dfrac{\partial\mathcal{L}}{\partial\left(\dfrac{\partial f}{\partial y}\right)}\right)
- \dfrac{\partial}{\partial z}
\left(\dfrac{\partial\mathcal{L}}{\partial\left(\dfrac{\partial f}{\partial z}\right)}\right)
= 0.
\end{aligned}
\]

La estructura es la misma que en el caso anterior: además del término
temporal, aparece un término de Euler-Lagrange por cada dirección
espacial. Para varios campos \(f_a\), se obtiene una ecuación de esta
forma para cada índice \(a\).

\subsubsection{3.2. Ejemplo: la
braquistócrona}\label{ejemplo-la-braquistuxf3crona}

Como ejemplo de optimización de un funcional fuera de la mecánica
analítica, se considera el problema de la braquistócrona, planteado por
Johann Bernoulli en 1696 y considerado uno de los orígenes del cálculo
de variaciones.

Sean dos puntos: \(A\), situado en el origen \((0,0)\), y
\(B = (x_B, y_B)\), situado más abajo y desplazado horizontalmente. Se
busca la forma de un alambre, descrito por una curva \(y(x)\) que une
\(A\) con \(B\), tal que una cuenta que parte del reposo en \(A\) y
desliza sin rozamiento por él bajo la acción de la gravedad llegue a
\(B\) en el menor tiempo posible. Para simplificar los signos, el eje
\(+y\) se orienta hacia abajo, en el sentido de la gravedad.

La magnitud que se minimiza es el tiempo total de recorrido \(T\), que
es un funcional de \(y(x)\). Se construye en tres pasos:

\begin{enumerate}
\def\labelenumi{\arabic{enumi}.}
\tightlist
\item
  En un tramo de longitud \(ds\) recorrido con rapidez \(v\), el tiempo
  empleado es \(dt = \dfrac{ds}{v}\).
\item
  Por el teorema de Pitágoras, el elemento de arco es
  \(ds = \sqrt{dx^2 + dy^2} = \sqrt{1 + (y')^2} \cdot dx\), donde
  \(y' = \dfrac{dy}{dx}\).
\item
  Por conservación de la energía mecánica, y puesto que la cuenta parte
  del reposo en \(y=0\), la energía cinética ganada es igual a la
  energía potencial perdida tras descender una altura \(y\):
  \(\dfrac{1}{2} \cdot m \cdot v^2 = m \cdot g \cdot y\), de donde
  \(v = \sqrt{2 \cdot g \cdot y}\). La masa se cancela, así que la curva
  óptima no depende de ella.
\end{enumerate}

Sumando las contribuciones \(dt\) a lo largo de la curva, el funcional
del tiempo queda:

\[
\begin{aligned}
T[y(x)] = \int_{0}^{x_B} \dfrac{\sqrt{1 + (y')^2}}{\sqrt{2 \cdot g \cdot y}} \, dx
\end{aligned}
\]

El integrando desempeña aquí el papel de la función \(F(x,y,y')\) del
apartado 3.1.2. Aunque este problema no procede de la mecánica
analítica, se suele denotar con la misma letra que un lagrangiano,
porque las ecuaciones que se le aplican son las mismas:

\[
\begin{aligned}
\mathcal{L}(y, y', x) = \dfrac{\sqrt{1 + (y')^2}}{\sqrt{2 \cdot g \cdot y}}
\end{aligned}
\]

\paragraph{Resolución mediante la identidad de
Beltrami}\label{resoluciuxf3n-mediante-la-identidad-de-beltrami}

La curva \(y(x)\) que minimiza este funcional debe satisfacer la
ecuación de Euler-Lagrange unidimensional (la del apartado 3.1.2,
multiplicada por \(-1\)):

\[
\begin{aligned}
\dfrac{d}{dx}\left( \dfrac{\partial \mathcal{L}}{\partial y'} \right) - \dfrac{\partial \mathcal{L}}{\partial y} = 0
\end{aligned}
\]

Aplicarla directamente exige calcular
\(\dfrac{d}{dx}\left( \dfrac{\partial \mathcal{L}}{\partial y'} \right)\),
lo que produce una ecuación de segundo orden con muchos términos. Se
puede evitar observando que \(\mathcal{L}\) no depende explícitamente de
\(x\): en su expresión solo aparecen \(y\) e \(y'\).

En ese caso, la ecuación de Euler-Lagrange admite una primera integral
(una cantidad que se conserva a lo largo de la solución), conocida como
identidad de Beltrami:

\[
\begin{aligned}
\mathcal{L} - y' \cdot \dfrac{\partial \mathcal{L}}{\partial y'} = C
\end{aligned}
\]

donde \(C\) es una constante. Para justificarla, se deriva el lado
izquierdo respecto de \(x\). Como \(\mathcal{L}\) depende de \(x\) a
través de \(y\) y de \(y'\) (y, en general, también explícitamente), la
regla de la cadena da:

\[
\begin{aligned}
\dfrac{d}{dx}\left( \mathcal{L} - y' \cdot \dfrac{\partial \mathcal{L}}{\partial y'} \right)
&= \dfrac{\partial \mathcal{L}}{\partial x} + y' \cdot \dfrac{\partial \mathcal{L}}{\partial y} + y'' \cdot \dfrac{\partial \mathcal{L}}{\partial y'}
- y'' \cdot \dfrac{\partial \mathcal{L}}{\partial y'} - y' \cdot \dfrac{d}{dx}\left( \dfrac{\partial \mathcal{L}}{\partial y'} \right) \\
&= \dfrac{\partial \mathcal{L}}{\partial x} - y' \cdot \left[ \dfrac{d}{dx}\left( \dfrac{\partial \mathcal{L}}{\partial y'} \right) - \dfrac{\partial \mathcal{L}}{\partial y} \right].
\end{aligned}
\]

El corchete es nulo cuando \(y(x)\) satisface la ecuación de
Euler-Lagrange, y \(\dfrac{\partial \mathcal{L}}{\partial x} = 0\) por
hipótesis. Por tanto, la derivada es cero y la expresión es constante.
La identidad sustituye una ecuación de segundo orden por otra de primer
orden.

Se aplica a continuación. Primero se calcula la derivada parcial de
\(\mathcal{L}\) respecto a \(y'\), tratando \(y\) como constante:

\[
\begin{aligned}
\dfrac{\partial \mathcal{L}}{\partial y'} = \dfrac{1}{\sqrt{2 \cdot g \cdot y}} \cdot \dfrac{1}{2 \cdot \sqrt{1 + (y')^2}} \cdot 2 \cdot y' = \dfrac{y'}{\sqrt{2 \cdot g \cdot y} \cdot \sqrt{1 + (y')^2}}
\end{aligned}
\]

Ahora introducimos \(\mathcal{L}\) y esta derivada en la Identidad de
Beltrami:

\[
\begin{aligned}
\dfrac{\sqrt{1 + (y')^2}}{\sqrt{2 \cdot g \cdot y}} - y' \cdot \left( \dfrac{y'}{\sqrt{2 \cdot g \cdot y} \cdot \sqrt{1 + (y')^2}} \right) = C
\end{aligned}
\]

Ambos términos del lado izquierdo llevan el factor
\(\dfrac{1}{\sqrt{2 \cdot g}}\). Multiplicando la ecuación por
\(\sqrt{2 \cdot g}\) y definiendo la nueva constante
\(C_1 = C \cdot \sqrt{2 \cdot g}\), queda:

\[
\begin{aligned}
\dfrac{\sqrt{1 + (y')^2}}{\sqrt{y}} - \dfrac{(y')^2}{\sqrt{y} \cdot \sqrt{1 + (y')^2}} = C_1
\end{aligned}
\]

Para restar estas fracciones, multiplicamos el primer término arriba y
abajo por \(\sqrt{1 + (y')^2}\) para buscar un denominador común:

\[
\begin{aligned}
\dfrac{1 + (y')^2 - (y')^2}{\sqrt{y} \cdot \sqrt{1 + (y')^2}} = C_1 \implies \dfrac{1}{\sqrt{y \cdot (1 + (y')^2)}} = C_1
\end{aligned}
\]

Elevando al cuadrado ambos lados e invirtiendo, se obtiene de nuevo una
constante, \(1/C_1^2\). Se la denota \(2 \cdot r\) porque, como se verá
más abajo, \(r\) resulta ser el radio de la circunferencia que genera la
curva solución:

\[
\begin{aligned}
y \cdot (1 + (y')^2) = \dfrac{1}{C_1^2} = 2 \cdot r
\end{aligned}
\]

Esta es la ecuación diferencial de primer orden que satisface la curva
óptima. Despejando la pendiente \(y' = \dfrac{dy}{dx}\) y tomando la
raíz positiva (en el tramo inicial la cuenta desciende, de modo que
\(y\) crece con \(x\)):

\[
\begin{aligned}
(y')^2 = \dfrac{2 \cdot r - y}{y} \implies \dfrac{dy}{dx} = \sqrt{\dfrac{2 \cdot r - y}{y}}
\end{aligned}
\]

Esta es la ecuación diferencial de una cicloide. Con el eje \(y\)
orientado hacia abajo, la curva tiene el aspecto de una cicloide
invertida respecto a su representación habitual. Su solución en forma
paramétrica, con parámetro \(\theta\) y condición inicial
\((x,y) = (0,0)\) en \(\theta = 0\), es:

\[
\begin{aligned}
x(\theta) = r \cdot (\theta - \sin \theta)
\end{aligned}
\]

\[
\begin{aligned}
y(\theta) = r \cdot (1 - \cos \theta)
\end{aligned}
\]

Comprobación: derivando respecto de \(\theta\),
\(\dfrac{dx}{d\theta} = r \cdot (1 - \cos\theta)\) y
\(\dfrac{dy}{d\theta} = r \cdot \sin\theta\), luego

\[
\begin{aligned}
\left(\dfrac{dy}{dx}\right)^2
= \dfrac{\sin^2\theta}{(1-\cos\theta)^2}
= \dfrac{(1-\cos\theta) \cdot (1+\cos\theta)}{(1-\cos\theta)^2}
= \dfrac{1+\cos\theta}{1-\cos\theta}.
\end{aligned}
\]

Por otro lado, sustituyendo \(y = r \cdot (1-\cos\theta)\) en el lado
derecho de la ecuación diferencial,
\(\dfrac{2 \cdot r - y}{y} = \dfrac{r \cdot (1+\cos\theta)}{r \cdot (1-\cos\theta)} = \dfrac{1+\cos\theta}{1-\cos\theta}\).
Ambas expresiones coinciden.

Las dos constantes que quedan, \(r\) y el valor final del parámetro
\(\theta_B\), se determinan imponiendo que la curva pase por \(B\):
\(x_B = r \cdot (\theta_B - \sin\theta_B)\) e
\(y_B = r \cdot (1 - \cos\theta_B)\).

Geométricamente, una cicloide es la curva que describe un punto del
borde de una rueda de radio \(r\) cuando esta rueda sin deslizar sobre
una recta; \(\theta\) es el ángulo girado por la rueda. Desde el punto
de vista físico, el resultado indica que la trayectoria de tiempo mínimo
no es la de longitud mínima (el segmento recto \(AB\)): la cicloide
empieza con una pendiente más pronunciada (vertical en \(A\), pues
\(dy/dx \to \infty\) cuando \(y \to 0\)), de modo que la cuenta adquiere
velocidad al principio del recorrido, y esa mayor velocidad compensa la
mayor longitud del camino.

(Nota: otro problema clásico del mismo tipo es el problema
isoperimétrico: entre todas las curvas cerradas de longitud fija,
encontrar la que encierra el área máxima. La solución es la
circunferencia. A diferencia de la braquistócrona, este problema incluye
una restricción (la longitud fija), que se incorpora mediante un
multiplicador de Lagrange).

\subsubsection{3.3. Variaciones y derivada
funcional}\label{variaciones-y-derivada-funcional}

Para estudiar si un funcional es estacionario, se compara una función
candidata \(f_0\) con funciones próximas de la forma
\(f_\epsilon = f_0 + \epsilon \cdot \eta\), donde \(\epsilon\) es un
parámetro real pequeño y \(\eta\) es una perturbación suave que respeta
las condiciones de contorno. La \textbf{primera variación} es el cambio
lineal del funcional al variar \(\epsilon\) alrededor de cero.

La derivada funcional \(\dfrac{\delta J}{\delta f}\) desempeña para un
funcional \(J[f]\) el papel que la derivada ordinaria desempeña para una
función de una variable. La condición \(\delta J = 0\) expresa que
ninguna perturbación admisible cambia el valor de \(J\) en primer orden.
En el caso de la acción mecánica, esta condición conduce a las
ecuaciones de Euler-Lagrange.

\subsection{4. El formalismo
lagrangiano}\label{el-formalismo-lagrangiano}

El formalismo lagrangiano es una reformulación de la mecánica clásica
que, en lugar de basarse en fuerzas vectoriales, se fundamenta en
principios energéticos y variacionales. Trabaja con magnitudes escalares
(energías) en lugar de vectoriales (fuerzas), y se extiende de forma
directa a medios continuos y campos, como los que describen la
propagación de ondas.

Para un sistema continuo, por ejemplo el desplazamiento transversal de
una cuerda, el estado está descrito por un campo \(f(x,t)\). La densidad
lagrangiana \(\mathcal{L}\) asigna una expresión local a cada punto del
espacio-tiempo a partir del campo, de sus derivadas y, en el caso
general, de la posición y del tiempo de forma explícita:

\[
\begin{aligned}
\mathcal{L}
= \mathcal{L}\left(
f,
\dfrac{\partial f}{\partial t},
\dfrac{\partial f}{\partial x},
x,t
\right).
\end{aligned}
\]

La dependencia explícita de \(x\) o \(t\) no debe confundirse con la
dependencia implícita que ya aparece a través de \(f(x,t)\). Si
\(\mathcal{L}\) no depende explícitamente de \(x\), el sistema es
homogéneo en el espacio; si no depende explícitamente de \(t\), es
homogéneo en el tiempo. Esta última simetría está asociada, por el
teorema de Noether, a la conservación de la energía. Una cuerda cuyo
tensado cambia externamente con el tiempo es un ejemplo de sistema con
dependencia explícita de \(t\).

La acción \(S\) es el funcional que se obtiene al integrar esa densidad
sobre la región de espacio-tiempo considerada:

\[
\begin{aligned}
S[f]
= \int_{t_1}^{t_2} \int_{x_1}^{x_2}
\mathcal{L}\left(
f,
\dfrac{\partial f}{\partial t},
\dfrac{\partial f}{\partial x},
x,t
\right) \, dx \, dt.
\end{aligned}
\]

El Principio de Hamilton establece que el campo físico que realmente se
da es aquel para el que esta acción es estacionaria: \(\delta S = 0\).

\subsubsection{4.1. Ecuaciones de Euler-Lagrange para campos
continuos}\label{ecuaciones-de-euler-lagrange-para-campos-continuos}

Volviendo al sistema físico continuo, para encontrar el campo \(f\) que
hace estacionaria la acción se somete el campo a una variación
infinitesimal \(f \rightarrow f + \delta f\). Asumimos que esta
variación se anula en los bordes de integración espaciales y temporales
(es decir, el sistema empieza y acaba en estados fijos conocidos, por lo
que \(\delta f = 0\) en \(x_1, x_2, t_1, t_2\)).

La variación total de la Acción \(S\) ante este pequeño cambio
arbitrario de la función viene dada por la regla de la cadena para
múltiples variables, aplicada al integrando \(\mathcal{L}\):

\[
\begin{aligned}
\delta S = \int \int \left[ \dfrac{\partial \mathcal{L}}{\partial f} \cdot \delta f + \dfrac{\partial \mathcal{L}}{\partial \left( \dfrac{\partial f}{\partial t} \right)} \cdot \delta \left( \dfrac{\partial f}{\partial t} \right) + \dfrac{\partial \mathcal{L}}{\partial \left( \dfrac{\partial f}{\partial x} \right)} \cdot \delta \left( \dfrac{\partial f}{\partial x} \right) \right] \, dx \, dt = 0
\end{aligned}
\]

\paragraph{\texorpdfstring{Justificación: conmutación de \(\delta\) y
\(\partial\)}{Justificación: conmutación de \textbackslash delta y \textbackslash partial}}\label{justificaciuxf3n-conmutaciuxf3n-de-delta-y-partial}

Para continuar es necesario intercambiar el operador variación
(\(\delta\)) con la derivada espacial o temporal, es decir, usar que
\(\delta(\partial_t f) = \partial_t (\delta f)\) y
\(\delta(\partial_x f) = \partial_x (\delta f)\).

Para justificarlo hay que precisar qué es una variación \(\delta f\).
Sea \(f_0(x,t)\) el campo físico (la solución buscada). Se perturba
sumándole una función suave arbitraria \(\eta(x,t)\), que se anula en
las fronteras, multiplicada por un parámetro real pequeño \(\epsilon\):

\[
\begin{aligned}
f(x,t,\epsilon) = f_0(x,t) + \epsilon \cdot \eta(x,t)
\end{aligned}
\]

El operador variación \(\delta\) representa matemáticamente la
diferencial respecto a este nuevo parámetro \(\epsilon\), evaluada en
cero:

\[
\begin{aligned}
\delta f = \left. \dfrac{\partial f}{\partial \epsilon} \right\vert{}_{\epsilon=0} \cdot d\epsilon = \eta(x,t) \cdot d\epsilon
\end{aligned}
\]

La posición \(x\), el tiempo \(t\) y el parámetro \(\epsilon\) son
variables independientes de la función \(f(x,t,\epsilon)\). Si \(f_0\) y
\(\eta\) son de clase \(C^2\), el teorema de Clairaut-Schwarz garantiza
que las derivadas parciales mixtas respecto de \(\epsilon\) y de \(t\)
(o de \(x\)) conmutan. Por lo tanto:

\[
\begin{aligned}
\delta \left( \dfrac{\partial f}{\partial t} \right) = \dfrac{\partial}{\partial \epsilon} \left( \dfrac{\partial f}{\partial t} \right) \cdot d\epsilon = \dfrac{\partial}{\partial t} \left( \dfrac{\partial f}{\partial \epsilon} \right) \cdot d\epsilon = \dfrac{\partial}{\partial t} (\delta f)
\end{aligned}
\]

El mismo argumento, con \(x\) en lugar de \(t\), da
\(\delta(\partial_x f) = \partial_x(\delta f)\). En resumen, como
\(\epsilon\) es independiente de \(x\) y de \(t\), variar y derivar son
operaciones que conmutan.

\paragraph{Integración por partes y derivada
funcional}\label{integraciuxf3n-por-partes-y-derivada-funcional}

Volviendo a la ecuación \(\delta S=0\), se usa la conmutación para
escribir el segundo y el tercer término con derivadas de \(\delta f\), y
después se integran por partes (la regla de la derivada del producto,
leída en sentido inverso). Para el término temporal, con \(x\) fijo:

\[
\begin{aligned}
\int_{t_1}^{t_2} \dfrac{\partial \mathcal{L}}{\partial \left(\dfrac{\partial f}{\partial t}\right)} \cdot \dfrac{\partial (\delta f)}{\partial t} \, dt
= \left[ \dfrac{\partial \mathcal{L}}{\partial \left(\dfrac{\partial f}{\partial t}\right)} \cdot \delta f \right]_{t_1}^{t_2}
- \int_{t_1}^{t_2} \dfrac{\partial}{\partial t}\left( \dfrac{\partial \mathcal{L}}{\partial \left(\dfrac{\partial f}{\partial t}\right)} \right) \cdot \delta f \, dt
\end{aligned}
\]

y análogamente para el término espacial, integrando en \(x\) con \(t\)
fijo.

Los términos de borde (los corchetes evaluados en \(t_1, t_2\) o en
\(x_1, x_2\)) se anulan porque, por hipótesis, la variación es nula en
los bordes (\(\delta f = 0\)). Solo quedan las integrales, con signo
menos:

\[
\begin{aligned}
\delta S = \int \int \left[ \dfrac{\partial \mathcal{L}}{\partial f} - \dfrac{\partial}{\partial t}\left( \dfrac{\partial \mathcal{L}}{\partial \left(\dfrac{\partial f}{\partial t}\right)} \right) - \dfrac{\partial}{\partial x}\left( \dfrac{\partial \mathcal{L}}{\partial \left(\dfrac{\partial f}{\partial x}\right)} \right) \right] \cdot \delta f \, dx \, dt = 0
\end{aligned}
\]

Al corchete de esta integral se le conoce matemáticamente como la
derivada funcional de la Acción \(S\) respecto al campo \(f(x,t)\)
(denotada comúnmente como \(\dfrac{\delta S}{\delta f}\)).

Para que esta integral doble sea cero para cualquier elección de la
variación \(\delta f\), el corchete (la derivada funcional) debe
anularse en todo punto del espacio-tiempo. Este es el lema fundamental
del cálculo de variaciones. (Idea de la demostración: si el corchete
fuera positivo en algún punto, por continuidad lo sería en un entorno;
eligiendo \(\delta f\) positiva en ese entorno y nula fuera, la integral
sería positiva, en contradicción con \(\delta S = 0\).)

Se obtiene así la ecuación de Euler-Lagrange para campos continuos,
escrita aquí con el signo cambiado respecto a la del apartado 3.1.3
(ambas formas son equivalentes):

\[
\begin{aligned}
\dfrac{\partial}{\partial t}\left( \dfrac{\partial \mathcal{L}}{\partial \left(\dfrac{\partial f}{\partial t}\right)} \right) + \dfrac{\partial}{\partial x}\left( \dfrac{\partial \mathcal{L}}{\partial \left(\dfrac{\partial f}{\partial x}\right)} \right) - \dfrac{\partial \mathcal{L}}{\partial f} = 0
\end{aligned}
\]

Todo sistema físico continuo descrito por una densidad lagrangiana
\(\mathcal{L}\) obedece esta ecuación diferencial, que es la condición
necesaria para que su acción sea estacionaria. El principio de Hamilton
exige estacionariedad, no necesariamente un mínimo: la solución puede
corresponder a un mínimo, a un máximo o a un punto de silla de la
acción.

\subsubsection{4.2. Sistemas discretos de
partículas}\label{sistemas-discretos-de-partuxedculas}

Hasta ahora hemos tratado medios continuos mediante una densidad
lagrangiana (\(\mathcal{L}\)) integrada sobre el volumen. Sin embargo,
en mecánica clásica es muy habitual tratar con sistemas discretos
compuestos por \(N\) puntos materiales (partículas puntuales) con masa.

Para este tipo de sistemas, ya no integramos sobre el espacio, sino que
definimos la posición del sistema mediante un conjunto de coordenadas
generalizadas \(q_i\). Si tenemos \(N\) partículas moviéndose libremente
en 3 dimensiones, necesitaremos \(3 \cdot N\) coordenadas (por ejemplo,
\(q_1 = x_1, q_2 = y_1, q_3 = z_1, q_4 = x_2...\) hasta \(q_{3N}\)). Si
existen \(r\) ligaduras independientes (por ejemplo, una distancia fija
entre dos partículas), el número de coordenadas independientes, o grados
de libertad, se reduce a \(n = 3 \cdot N - r\). Las coordenadas
generalizadas no tienen por qué ser cartesianas: pueden ser ángulos,
distancias u otras variables que describan la configuración.

Las derivadas temporales de estas coordenadas se denominan velocidades
generalizadas, expresadas mediante la notación de punto de Newton:

\[
\begin{aligned}
\dot{q}_i = \dfrac{dq_i}{dt}
\end{aligned}
\]

En este contexto discreto, la densidad lagrangiana se sustituye por una
función escalar llamada Lagrangiano (\(L\)), que depende explícitamente
de las coordenadas generalizadas, las velocidades generalizadas y,
opcionalmente, del tiempo:

\[
\begin{aligned}
L = L(q_1, q_2, ..., q_{n}, \dot{q}_1, \dot{q}_2, ..., \dot{q}_{n}, t) \equiv L(q_i, \dot{q}_i, t)
\end{aligned}
\]

donde \(n\) es el número de grados de libertad (\(n = 3 \cdot N\) si no
hay ligaduras).

La Acción se define ahora como una única integral en el tiempo del
Lagrangiano a lo largo de la trayectoria del sistema entre un estado
inicial \(t_1\) y uno final \(t_2\):

\[
\begin{aligned}
S = \int_{t_1}^{t_2} L(q_i, \dot{q}_i, t) \, dt
\end{aligned}
\]

Aplicando el Principio de Hamilton (\(\delta S = 0\)) y el mismo
procedimiento del apartado 3.1.2 (ahora con una sola variable de
integración, el tiempo, y una variación independiente \(\delta q_i\)
para cada coordenada), obtenemos un conjunto de \(n\) ecuaciones
diferenciales ordinarias acopladas, una por coordenada
(\(i = 1, \dots, n\)):

\[
\begin{aligned}
\dfrac{d}{dt}\left( \dfrac{\partial L}{\partial \dot{q}_i} \right) - \dfrac{\partial L}{\partial q_i} = 0
\end{aligned}
\]

Estas son las Ecuaciones de Euler-Lagrange para sistemas discretos. Cada
término tiene un significado físico directo:

\begin{enumerate}
\def\labelenumi{\arabic{enumi}.}
\tightlist
\item
  El término \(\dfrac{\partial L}{\partial \dot{q}_i}\) se define como
  el momento generalizado (o momento conjugado) \(p_i\). (Si
  \(L = \dfrac{1}{2} \cdot m \cdot \dot{x}^2 - V(x)\), entonces
  \(\dfrac{\partial L}{\partial \dot{x}} = m \cdot \dot{x} = p_x\)).
\item
  El término \(\dfrac{\partial L}{\partial q_i}\) se define como la
  fuerza generalizada \(F_i\). (Si \(V\) solo depende de \(x\), entonces
  \(\dfrac{\partial L}{\partial x} = -\dfrac{\partial V}{\partial x} = F_x\)).
\end{enumerate}

Sustituyendo estas definiciones, la ecuación se lee como
\(\dfrac{d p_i}{dt} = F_i\), que en coordenadas cartesianas es la
Segunda Ley de Newton. La diferencia respecto a la formulación
newtoniana es que las ecuaciones de Euler-Lagrange tienen la misma forma
en cualquier sistema de coordenadas (cartesianas, polares, cilíndricas o
variables abstractas). En coordenadas no cartesianas, la energía
cinética puede depender de \(q_i\) (por ejemplo, en polares
\(T = \dfrac{1}{2} \cdot m \cdot (\dot{r}^2 + r^2 \cdot \dot{\theta}^2)\)
depende de \(r\)), y entonces \(\dfrac{\partial L}{\partial q_i}\)
incluye, además de la fuerza, términos procedentes de \(T\), como la
fuerza centrífuga \(m \cdot r \cdot \dot{\theta}^2\) en la ecuación de
\(r\).

\subsubsection{\texorpdfstring{4.3. Por qué \(L = T - V\): el puente
desde
Newton}{4.3. Por qué L = T - V: el puente desde Newton}}\label{por-quuxe9-l-t---v-el-puente-desde-newton}

Es común ver definido el Lagrangiano como la diferencia entre la Energía
Cinética Total (\(T\)) y la Energía Potencial (\(V\)):

\[
\begin{aligned}
L = T - V
\end{aligned}
\]

Esta forma no es arbitraria: es la que hace que las ecuaciones de
Euler-Lagrange coincidan con la Segunda Ley de Newton cuando las fuerzas
son conservativas.

Se muestra a continuación para una única partícula en una dimensión
(coordenada \(x\)). La Segunda Ley de Newton nos dice:

\[
\begin{aligned}
m \cdot \ddot{x} = F_x
\end{aligned}
\]

Si la fuerza es conservativa, deriva de un potencial \(V(x)\), de modo
que \(F_x = -\dfrac{\partial V}{\partial x}\). Sustituyendo:

\[
\begin{aligned}
m \cdot \ddot{x} = -\dfrac{\partial V}{\partial x}
\end{aligned}
\]

Por otro lado, la energía cinética de la partícula es
\(T = \dfrac{1}{2} \cdot m \cdot \dot{x}^2\). Notemos dos propiedades
matemáticas de \(T\):

\begin{enumerate}
\def\labelenumi{\arabic{enumi}.}
\tightlist
\item
  Si derivamos \(T\) respecto a la velocidad \(\dot{x}\), obtenemos el
  momento: \(\dfrac{\partial T}{\partial \dot{x}} = m \cdot \dot{x}\).
  Si a esto le aplicamos la derivada temporal, recuperamos la fuerza
  inercial (masa por aceleración):
  \(\dfrac{d}{dt}\left( \dfrac{\partial T}{\partial \dot{x}} \right) = m \cdot \ddot{x}\).
\item
  Si derivamos \(T\) respecto a la posición \(x\), el resultado es cero
  (en coordenadas cartesianas, la energía cinética depende solo de la
  velocidad de la partícula, no de su posición):
  \(\dfrac{\partial T}{\partial x} = 0\).
\end{enumerate}

Con esto en mente, podemos reescribir el término \(m \cdot \ddot{x}\) de
la ley de Newton usando exclusivamente derivadas de la energía cinética:

\[
\begin{aligned}
\dfrac{d}{dt}\left( \dfrac{\partial T}{\partial \dot{x}} \right) = -\dfrac{\partial V}{\partial x}
\end{aligned}
\]

Reagrupemos todos los términos en el lado izquierdo:

\[
\begin{aligned}
\dfrac{d}{dt}\left( \dfrac{\partial T}{\partial \dot{x}} \right) + \dfrac{\partial V}{\partial x} = 0
\end{aligned}
\]

El último paso consiste en sumar términos nulos. Como la energía
potencial \(V(x)\) depende únicamente de la posición, su derivada
parcial respecto a la velocidad es cero:
\(\dfrac{\partial V}{\partial \dot{x}} = 0\). Y, como se vio en el punto
2, \(\dfrac{\partial T}{\partial x} = 0\). Por tanto:

\begin{itemize}
\tightlist
\item
  \(\dfrac{\partial (T - V)}{\partial \dot{x}} = \dfrac{\partial T}{\partial \dot{x}} - 0 = \dfrac{\partial T}{\partial \dot{x}}\).
\item
  \(-\dfrac{\partial (T - V)}{\partial x} = -0 + \dfrac{\partial V}{\partial x} = \dfrac{\partial V}{\partial x}\).
\end{itemize}

Sustituyendo ambas igualdades en la ecuación anterior, las dos energías
quedan agrupadas bajo las mismas derivadas sin que cambie la ecuación:

\[
\begin{aligned}
\dfrac{d}{dt}\left( \dfrac{\partial (T - V)}{\partial \dot{x}} \right) - \dfrac{\partial (T - V)}{\partial x} = 0
\end{aligned}
\]

Esta ecuación tiene la forma de la ecuación de Euler-Lagrange,
\(\dfrac{d}{dt}\left( \dfrac{\partial L}{\partial \dot{x}} \right) - \dfrac{\partial L}{\partial x} = 0\),
con \(L = T - V\). Es decir, tomando \(L = T - V\), el principio
variacional reproduce la dinámica newtoniana de una partícula sometida a
fuerzas conservativas.

Esta elección no es la única posible: multiplicar \(L\) por una
constante no nula, o sumarle la derivada temporal total de una función
\(G(q,t)\), produce las mismas ecuaciones de movimiento. \(L = T - V\)
es la elección estándar.

(Nota: En sistemas mecánicos con ligaduras y coordenadas generalizadas
\(q_i\), se utiliza el Principio de d'Alembert (o de los Trabajos
Virtuales) para demostrar de forma general que la misma estructura
\(L = T - V\) sigue siendo válida, siempre que las fuerzas deriven de un
potencial y las ligaduras no realicen trabajo).

\subsection{5. Introducción a la transformada de
Legendre}\label{introducciuxf3n-a-la-transformada-de-legendre}

La transformada de Legendre permite describir una misma función mediante
su pendiente, en lugar de mediante su variable original. No se trata de
renombrar la variable: la transformada sustituye la variable \(v\) por
la cantidad conjugada \(p\), definida como la pendiente de la gráfica de
\(F\):

\[
\begin{aligned}
p = \dfrac{dF}{dv}.
\end{aligned}
\]

\subsubsection{5.1. Motivación geométrica y
definición}\label{motivaciuxf3n-geomuxe9trica-y-definiciuxf3n}

Supongamos que \(F(v)\) es diferenciable y convexa. En cada valor de
\(v\), la recta tangente a su gráfica tiene pendiente \(p=F'(v)\) y
ecuación:

\[
\begin{aligned}
\ell(u) = F(v) + p \cdot (u-v).
\end{aligned}
\]

Si evaluamos esta recta en \(u=0\), obtenemos su ordenada en el origen,
\(\ell(0)=F(v)-p \cdot v\). La cantidad opuesta, \(p \cdot v-F(v)\),
depende de la pendiente \(p\) y permite codificar la misma familia de
rectas tangentes. Por definición, la transformada de Legendre de \(F\)
es:

\[
\begin{aligned}
G(p) = p \cdot v - F(v),
\end{aligned}
\]

donde \(v\) debe expresarse como función de \(p\) mediante \(p=F'(v)\).
La convexidad asegura localmente que esta relación puede invertirse y
que cada pendiente corresponde a un único punto de contacto. Por ello,
la transformación conserva la información de \(F\), pero la describe
mediante su variable conjugada \(p\).

Al derivar \(G(p)=p \cdot v(p)-F(v(p))\) respecto de \(p\), se ve el
intercambio de variables de manera algebraica:

\[
\begin{aligned}
\dfrac{dG}{dp}
&= v + p \cdot \dfrac{dv}{dp} - \dfrac{dF}{dv} \cdot \dfrac{dv}{dp} \\
&= v + \left(p-\dfrac{dF}{dv}\right) \cdot \dfrac{dv}{dp} \\
&= v.
\end{aligned}
\]

La última igualdad usa precisamente \(p=\dfrac{dF}{dv}\). La
transformada intercambia así los papeles de variable y pendiente: de
\(p=F'(v)\) se pasa a \(v=G'(p)\).

\subsubsection{5.2. Ejemplo cuadrático y doble
transformada}\label{ejemplo-cuadruxe1tico-y-doble-transformada}

Consideremos la función convexa:

\[
\begin{aligned}
F(v)=v^2.
\end{aligned}
\]

Su variable conjugada es la pendiente:

\[
\begin{aligned}
p = \dfrac{dF}{dv} = 2 \cdot v,
\qquad
v = \dfrac{p}{2}.
\end{aligned}
\]

Sustituimos esta expresión de \(v\) en la definición de la transformada:

\[
\begin{aligned}
G(p)
&= p \cdot v - F(v) \\
&= p \cdot \left(\dfrac{p}{2}\right) - \left(\dfrac{p}{2}\right)^2 \\
&= \dfrac{p^2}{4}.
\end{aligned}
\]

Podemos recuperar la función original aplicando de nuevo la transformada
de Legendre a \(G\). Ahora la variable conjugada de \(p\) es:

\[
\begin{aligned}
v = \dfrac{dG}{dp} = \dfrac{p}{2},
\qquad
p=2 \cdot v.
\end{aligned}
\]

La doble transformada vale entonces:

\[
\begin{aligned}
F^{**}(v)
&= v \cdot p-G(p) \\
&= v \cdot (2 \cdot v)-\dfrac{(2 \cdot v)^2}{4} \\
&= v^2 = F(v).
\end{aligned}
\]

Para funciones diferenciables, convexas y adecuadamente regulares, la
doble transformada devuelve la función original: \(F^{**}=F\). Esta
propiedad explica por qué la transformación de Legendre no pierde
información y por qué resulta adecuada para pasar del formalismo
lagrangiano al hamiltoniano.

\textbf{Ejemplo computacional:}
\href{https://github.com/julian1c2a/MCC/blob/main/cuadernos/tema1/legendre.py}{cuadernos/tema1/legendre.py}.
Comprueba con SymPy este ejemplo y estudia
\(f(x,t) = \cos(k \cdot x) \cdot e^{\omega \cdot t}\): transformada
respecto de \(x\) (ramas, variable pasiva, doble transformada) y
respecto de \(x\) y \(t\) a la vez.

\subsection{6. Introducción al formalismo
hamiltoniano}\label{introducciuxf3n-al-formalismo-hamiltoniano}

El formalismo lagrangiano \(L(q_i, \dot{q}_i, t)\) formula la dinámica
en el espacio de configuración, definido por las coordenadas
generalizadas \(q_i\). La mecánica hamiltoniana aplica la transformada
de Legendre a las velocidades \(\dot{q}_i\) para usar, en su lugar, los
momentos conjugados \(p_i\). Así se obtiene una descripción en el
espacio de fases, cuyas coordenadas son los pares \((q_i,p_i)\).

\subsubsection{6.1. Construcción del Hamiltoniano mediante la
transformada de
Legendre}\label{construcciuxf3n-del-hamiltoniano-mediante-la-transformada-de-legendre}

Como vimos en la sección anterior, el momento generalizado o momento
conjugado se define a partir del Lagrangiano como:

\[
\begin{aligned}
p_i = \dfrac{\partial L}{\partial \dot{q}_i}
\end{aligned}
\]

El objetivo es construir una nueva función, el Hamiltoniano (\(H\)), que
dependa exclusivamente de posiciones, momentos y del tiempo:
\(H(q_i, p_i, t)\). El cambio de variables \(\dot{q}_i \to p_i\) se
realiza mediante la transformada de Legendre de la sección 5.

La transformada de Legendre pasa de una función que depende de una
variable (en este caso, \(L\) como función de \(\dot{q}_i\)) a otra que
depende de la derivada parcial de la función original respecto a esa
variable (es decir, de \(p_i = \partial L/\partial \dot{q}_i\)). Las
variables \(q_i\) y \(t\) no intervienen en la transformación y actúan
como parámetros. Para que el cambio sea posible hay que poder despejar
las velocidades \(\dot{q}_i\) en función de \((q_i, p_i, t)\), lo que
exige que la matriz de segundas derivadas
\(\dfrac{\partial^2 L}{\partial \dot{q}_i \partial \dot{q}_j}\) sea
invertible (en el caso de una variable, \(L\) convexa en \(\dot{q}\),
como en la sección 5). La definición general del Hamiltoniano, análoga a
\(G(p) = p \cdot v - F(v)\) con una suma sobre todas las coordenadas,
es:

\[
\begin{aligned}
H(q_i, p_i, t) = \sum_{i} p_i \cdot \dot{q}_i - L(q_i, \dot{q}_i, t)
\end{aligned}
\]

\subsubsection{6.2. Deducción de las ecuaciones canónicas de
Hamilton}\label{deducciuxf3n-de-las-ecuaciones-canuxf3nicas-de-hamilton}

Para descubrir qué forma toman las ecuaciones de movimiento en este
nuevo marco, vamos a calcular la diferencial total (el cambio
infinitesimal) del Hamiltoniano de dos maneras distintas y compararlas.

\paragraph{Primer método: usando la definición de la transformada de
Legendre}\label{primer-muxe9todo-usando-la-definiciuxf3n-de-la-transformada-de-legendre}

Si diferenciamos la ecuación \(H = \sum p_i \cdot \dot{q}_i - L\),
aplicando la regla del producto y la regla de la cadena para la
diferencial de \(L(q_i, \dot{q}_i, t)\), obtenemos:

\[
\begin{aligned}
dH = \sum_{i} (\dot{q}_i \cdot dp_i + p_i \cdot d\dot{q}_i) - \left[ \sum_{i} \left( \dfrac{\partial L}{\partial q_i} \cdot dq_i + \dfrac{\partial L}{\partial \dot{q}_i} \cdot d\dot{q}_i \right) + \dfrac{\partial L}{\partial t} \cdot dt \right]
\end{aligned}
\]

Ahora usamos dos relaciones del formalismo lagrangiano:

\begin{enumerate}
\def\labelenumi{\arabic{enumi}.}
\tightlist
\item
  \(p_i = \dfrac{\partial L}{\partial \dot{q}_i}\) (definición de
  momento).
\item
  \(\dot{p}_i = \dfrac{\partial L}{\partial q_i}\) (es la propia
  Ecuación de Euler-Lagrange,
  \(\dfrac{d p_i}{dt} = \dfrac{\partial L}{\partial q_i}\)).
\end{enumerate}

Sustituyendo esto en el corchete:

\[
\begin{aligned}
dH = \sum_{i} (\dot{q}_i \cdot dp_i + p_i \cdot d\dot{q}_i) - \sum_{i} (\dot{p}_i \cdot dq_i + p_i \cdot d\dot{q}_i) - \dfrac{\partial L}{\partial t} \cdot dt
\end{aligned}
\]

El término \(p_i \cdot d\dot{q}_i\) aparece una vez sumando y otra
restando, por lo que se cancela. Esta cancelación es la propiedad
central de la transformada de Legendre: \(dH\) no contiene los
diferenciales de velocidad \(d\dot{q}_i\), lo que confirma que \(H\) no
depende de las velocidades, sino de \(q_i\), \(p_i\) y \(t\). Nos queda:

\[
\begin{aligned}
dH = \sum_{i} (\dot{q}_i \cdot dp_i - \dot{p}_i \cdot dq_i) - \dfrac{\partial L}{\partial t} \cdot dt
\end{aligned}
\]

\paragraph{Segundo método: diferenciando el Hamiltoniano como función de
sus propias
variables}\label{segundo-muxe9todo-diferenciando-el-hamiltoniano-como-funciuxf3n-de-sus-propias-variables}

Sabemos por definición que \(H\) es una función de \(q_i\), \(p_i\) y
\(t\). Su diferencial total formal en cálculo multivariable es:

\[
\begin{aligned}
dH = \sum_{i} \left( \dfrac{\partial H}{\partial q_i} \cdot dq_i + \dfrac{\partial H}{\partial p_i} \cdot dp_i \right) + \dfrac{\partial H}{\partial t} \cdot dt
\end{aligned}
\]

\paragraph{Igualación de
coeficientes}\label{igualaciuxf3n-de-coeficientes}

Ambos métodos calculan el mismo diferencial \(dH\). Como \(q_i\),
\(p_i\) y \(t\) son variables independientes en el espacio de fases, sus
diferenciales \(dq_i\), \(dp_i\) y \(dt\) también lo son, y los
coeficientes que los acompañan deben coincidir término a término:

\begin{itemize}
\tightlist
\item
  Coeficiente de \(dp_i\):
  \(\dfrac{\partial H}{\partial p_i} = \dot{q}_i\).
\item
  Coeficiente de \(dq_i\):
  \(\dfrac{\partial H}{\partial q_i} = -\dot{p}_i\).
\item
  Coeficiente de \(dt\):
  \(\dfrac{\partial H}{\partial t} = -\dfrac{\partial L}{\partial t}\).
\end{itemize}

Las dos primeras son las ecuaciones de movimiento de la mecánica
hamiltoniana, conocidas como Ecuaciones Canónicas de Hamilton:

\[
\begin{aligned}
\dot{q}_i = \dfrac{\partial H}{\partial p_i}
\end{aligned}
\]

\[
\begin{aligned}
\dot{p}_i = -\dfrac{\partial H}{\partial q_i}
\end{aligned}
\]

La tercera relaciona la dependencia temporal explícita de ambas
funciones: \(H\) depende explícitamente del tiempo si y solo si \(L\) lo
hace.

Con \(n\) grados de libertad, el formalismo lagrangiano da \(n\)
ecuaciones diferenciales de segundo orden en las \(q_i\), mientras que
el formalismo hamiltoniano da \(2 \cdot n\) ecuaciones diferenciales
acopladas de primer orden en las \(q_i\) y las \(p_i\). El número total
de condiciones iniciales necesarias es el mismo en ambos casos
(\(2 \cdot n\): posiciones y velocidades, o posiciones y momentos). Dado
un punto inicial \((q_i(t_0), p_i(t_0))\), el teorema de existencia y
unicidad de las ecuaciones diferenciales ordinarias garantiza (para
\(H\) suficientemente regular) que existe una única trayectoria en el
espacio de fases, de dimensión \(2 \cdot n\), que pasa por él. En
consecuencia, dos trayectorias distintas del espacio de fases no se
cortan.

\subsection{7. Significado físico del
Hamiltoniano}\label{significado-fuxedsico-del-hamiltoniano}

El Hamiltoniano se ha definido como una transformada de Legendre del
Lagrangiano, \(H = \sum_i p_i \cdot \dot{q}_i - L\), sin ninguna
referencia a la energía. Esta sección estudia cuándo coincide con la
energía mecánica total \(E = T + V\) y cuándo se conserva. Son dos
preguntas distintas y, como se verá, sus respuestas son independientes.

\subsubsection{7.1. Funciones homogéneas y teorema de
Euler}\label{funciones-homoguxe9neas-y-teorema-de-euler}

Una función \(f(x_1, \dots, x_n)\) es \textbf{homogénea de grado \(k\)}
si, para todo \(\lambda > 0\),

\[
\begin{aligned}
f(\lambda \cdot x_1, \dots, \lambda \cdot x_n) = \lambda^k \cdot f(x_1, \dots, x_n).
\end{aligned}
\]

Por ejemplo, \(a \cdot x^2 + b \cdot x \cdot y + c \cdot y^2\) es
homogénea de grado 2, \(a \cdot x + b \cdot y\) lo es de grado 1 y una
función que no depende de las \(x_i\) lo es de grado 0.

\textbf{Teorema de Euler.} Si \(f\) es diferenciable y homogénea de
grado \(k\), entonces

\[
\begin{aligned}
\sum_{i=1}^{n} x_i \cdot \dfrac{\partial f}{\partial x_i} = k \cdot f.
\end{aligned}
\]

Demostración: se deriva la igualdad de la definición respecto de
\(\lambda\). En el lado izquierdo, por la regla de la cadena, cada
argumento \(\lambda \cdot x_i\) aporta su derivada \(x_i\):

\[
\begin{aligned}
\sum_{i=1}^{n} x_i \cdot \dfrac{\partial f}{\partial x_i}(\lambda \cdot x_1, \dots, \lambda \cdot x_n) = k \cdot \lambda^{k-1} \cdot f(x_1, \dots, x_n).
\end{aligned}
\]

Evaluando en \(\lambda = 1\) se obtiene el resultado.

En mecánica, la homogeneidad que interesa es respecto de las
\textbf{velocidades generalizadas} \(\dot{q}_i\), con las coordenadas
\(q_i\) y el tiempo \(t\) fijos. No se trata de homogeneidad en las
variables espaciales: \(T = \dfrac{1}{2} \cdot m \cdot \dot{x}^2\) es
homogénea de grado 2 en \(\dot{x}\), sea cual sea la dependencia en
\(x\).

\subsubsection{7.2. Descomposición de la energía
cinética}\label{descomposiciuxf3n-de-la-energuxeda-cinuxe9tica}

Sea un sistema de partículas de masas \(m_a\) cuyas posiciones se
expresan en función de las coordenadas generalizadas y, posiblemente,
del tiempo: \(\mathbf{r}_a = \mathbf{r}_a(q_1, \dots, q_n, t)\). La
dependencia explícita de \(t\) aparece cuando las ligaduras se mueven
(un alambre que gira, un soporte que oscila); en ese caso el sistema se
llama \textbf{reónomo}. Si \(\mathbf{r}_a\) no depende explícitamente de
\(t\), se llama \textbf{esclerónomo}.

Por la regla de la cadena, la velocidad de cada partícula es

\[
\begin{aligned}
\dot{\mathbf{r}}_a = \sum_{j} \dfrac{\partial \mathbf{r}_a}{\partial q_j} \cdot \dot{q}_j + \dfrac{\partial \mathbf{r}_a}{\partial t}.
\end{aligned}
\]

Al sustituir en
\(T = \dfrac{1}{2} \cdot \sum_a m_a \cdot \dot{\mathbf{r}}_a \cdot \dot{\mathbf{r}}_a\)
y desarrollar el cuadrado, la energía cinética se separa en tres partes,
homogéneas de grados 2, 1 y 0 en las velocidades:

\[
\begin{aligned}
T &= T_2 + T_1 + T_0, \\
T_2 &= \dfrac{1}{2} \cdot \sum_{j,k} M_{jk} \cdot \dot{q}_j \cdot \dot{q}_k,
\qquad M_{jk} = \sum_a m_a \cdot \dfrac{\partial \mathbf{r}_a}{\partial q_j} \cdot \dfrac{\partial \mathbf{r}_a}{\partial q_k}, \\
T_1 &= \sum_{j} b_j \cdot \dot{q}_j,
\qquad b_j = \sum_a m_a \cdot \dfrac{\partial \mathbf{r}_a}{\partial q_j} \cdot \dfrac{\partial \mathbf{r}_a}{\partial t}, \\
T_0 &= \dfrac{1}{2} \cdot \sum_a m_a \cdot \dfrac{\partial \mathbf{r}_a}{\partial t} \cdot \dfrac{\partial \mathbf{r}_a}{\partial t}.
\end{aligned}
\]

Los coeficientes \(M_{jk}\), \(b_j\) y \(T_0\) dependen de \(q\) y de
\(t\), pero no de las velocidades. En un sistema esclerónomo,
\(\dfrac{\partial \mathbf{r}_a}{\partial t} = \mathbf{0}\), de modo que
\(T_1 = T_0 = 0\) y \(T = T_2\) es homogénea de grado 2.

\subsubsection{7.3. Expresión general del
Hamiltoniano}\label{expresiuxf3n-general-del-hamiltoniano}

Del mismo modo, el Lagrangiano se descompone en partes homogéneas en las
velocidades, \(L = L_2 + L_1 + L_0\). Si el potencial \(V(q,t)\) no
depende de las velocidades, \(L_2 = T_2\), \(L_1 = T_1\) y
\(L_0 = T_0 - V\).

Como \(p_i = \dfrac{\partial L}{\partial \dot{q}_i}\), el teorema de
Euler aplicado a cada parte da

\[
\begin{aligned}
\sum_i p_i \cdot \dot{q}_i = \sum_i \dot{q}_i \cdot \dfrac{\partial L}{\partial \dot{q}_i} = 2 \cdot L_2 + 1 \cdot L_1 + 0 \cdot L_0,
\end{aligned}
\]

y, restando \(L\),

\[
\begin{aligned}
H = 2 \cdot L_2 + L_1 - (L_2 + L_1 + L_0) = L_2 - L_0.
\end{aligned}
\]

La parte lineal en las velocidades, \(L_1\), desaparece del
Hamiltoniano. Con \(V\) independiente de las velocidades:

\[
\begin{aligned}
H &= T_2 - T_0 + V, \\
E &= T + V = T_2 + T_1 + T_0 + V, \\
E - H &= T_1 + 2 \cdot T_0.
\end{aligned}
\]

\textbf{Condición para que \(H\) sea la energía total.} \(H = E\) para
todo estado del sistema si y solo si \(T_1 + 2 \cdot T_0 = 0\) para
cualesquiera velocidades. Como \(T_1\) es lineal en las velocidades y
\(T_0\) no depende de ellas, esto exige \(T_1 = 0\) y \(T_0 = 0\) por
separado; es decir, que \(T\) sea homogénea de grado 2 en las
velocidades. Esto se cumple siempre en los sistemas esclerónomos. Por
tanto, si \(V\) no depende de las velocidades:

\[
\begin{aligned}
H = T + V = E \iff T \text{ es homogénea de grado 2 en las } \dot{q}_i.
\end{aligned}
\]

El caso más sencillo es una partícula en coordenadas cartesianas, con
\(T = \dfrac{1}{2} \cdot m \cdot \dot{x}^2\) y \(V(x)\): entonces
\(p \cdot \dot{x} = \dot{x} \cdot \dfrac{\partial T}{\partial \dot{x}} = \dot{x} \cdot m \cdot \dot{x} = 2 \cdot T\),
y \(H = 2 \cdot T - (T - V) = T + V\).

\subsubsection{\texorpdfstring{7.4. Dependencia explícita del tiempo y
conservación de
\(H\)}{7.4. Dependencia explícita del tiempo y conservación de H}}\label{dependencia-expluxedcita-del-tiempo-y-conservaciuxf3n-de-h}

Usando la regla de la cadena y las ecuaciones canónicas:

\[
\begin{aligned}
\dfrac{dH}{dt}
&= \sum_{i} \left( \dfrac{\partial H}{\partial q_i} \cdot \dot{q}_i + \dfrac{\partial H}{\partial p_i} \cdot \dot{p}_i \right) + \dfrac{\partial H}{\partial t} \\
&= \sum_{i} \left( \dfrac{\partial H}{\partial q_i} \cdot \dfrac{\partial H}{\partial p_i} - \dfrac{\partial H}{\partial p_i} \cdot \dfrac{\partial H}{\partial q_i} \right) + \dfrac{\partial H}{\partial t} \\
&= \dfrac{\partial H}{\partial t} = -\dfrac{\partial L}{\partial t},
\end{aligned}
\]

donde la última igualdad es la tercera relación obtenida en la sección
6.2. Por tanto, \(H\) se conserva a lo largo del movimiento si y solo si
\(L\) (equivalentemente, \(H\)) no depende explícitamente del tiempo.

La conservación de \(H\) y su igualdad con \(E\) son propiedades
independientes:

\begin{itemize}
\tightlist
\item
  \(H = E\) depende de la forma de \(T\) (homogénea de grado 2) y de que
  \(V\) no dependa de las velocidades.
\item
  La conservación de \(H\) depende de que no haya dependencia explícita
  del tiempo en \(L\).
\end{itemize}

Un sistema reónomo puede tener un Lagrangiano sin dependencia explícita
del tiempo (por ejemplo, si la ligadura gira con velocidad angular
constante): entonces \(H\) se conserva aunque no sea la energía. En ese
caso, a la cantidad conservada \(H\) se la llama \textbf{integral de
Jacobi}. A la inversa, un sistema esclerónomo con un potencial que
depende del tiempo tiene \(H = E\), pero \(E\) no se conserva.

\subsubsection{7.5. Ejemplos}\label{ejemplos}

\paragraph{\texorpdfstring{Cuenta en una varilla que gira: \(H\) se
conserva, pero
\(H \neq E\)}{Cuenta en una varilla que gira: H se conserva, pero H \textbackslash neq E}}\label{cuenta-en-una-varilla-que-gira-h-se-conserva-pero-h-neq-e}

Una cuenta de masa \(m\) desliza sin rozamiento por una varilla recta
horizontal que gira en torno a un eje vertical con velocidad angular
constante \(\omega\). Se toma como coordenada generalizada la distancia
\(r\) de la cuenta al eje. Las coordenadas cartesianas son
\(x = r \cdot \cos(\omega \cdot t)\) e
\(y = r \cdot \sin(\omega \cdot t)\), que dependen explícitamente de
\(t\): el sistema es reónomo. Derivando:

\[
\begin{aligned}
\dot{x} &= \dot{r} \cdot \cos(\omega \cdot t) - r \cdot \omega \cdot \sin(\omega \cdot t), \\
\dot{y} &= \dot{r} \cdot \sin(\omega \cdot t) + r \cdot \omega \cdot \cos(\omega \cdot t), \\
T &= \dfrac{1}{2} \cdot m \cdot (\dot{x}^2 + \dot{y}^2) = \dfrac{1}{2} \cdot m \cdot \dot{r}^2 + \dfrac{1}{2} \cdot m \cdot \omega^2 \cdot r^2.
\end{aligned}
\]

Aquí \(T_2 = \dfrac{1}{2} \cdot m \cdot \dot{r}^2\), \(T_1 = 0\) y
\(T_0 = \dfrac{1}{2} \cdot m \cdot \omega^2 \cdot r^2\). La varilla es
horizontal, así que \(V = 0\) y \(L = T\). Entonces
\(p = \dfrac{\partial L}{\partial \dot{r}} = m \cdot \dot{r}\) y

\[
\begin{aligned}
H &= p \cdot \dot{r} - L = \dfrac{1}{2} \cdot m \cdot \dot{r}^2 - \dfrac{1}{2} \cdot m \cdot \omega^2 \cdot r^2 = T_2 - T_0, \\
E &= T = \dfrac{1}{2} \cdot m \cdot \dot{r}^2 + \dfrac{1}{2} \cdot m \cdot \omega^2 \cdot r^2.
\end{aligned}
\]

La diferencia es \(E - H = 2 \cdot T_0 = m \cdot \omega^2 \cdot r^2\),
como predice la sección 7.3.

\begin{itemize}
\tightlist
\item
  \(L\) no depende explícitamente de \(t\), luego \(H\) se conserva.
\item
  \(E\) no se conserva. La ecuación de Euler-Lagrange es
  \(m \cdot \ddot{r} - m \cdot \omega^2 \cdot r = 0\), es decir,
  \(\ddot{r} = \omega^2 \cdot r\), y entonces \[
  \begin{aligned}
  \dfrac{dE}{dt} = m \cdot \dot{r} \cdot \ddot{r} + m \cdot \omega^2 \cdot r \cdot \dot{r} = 2 \cdot m \cdot \omega^2 \cdot r \cdot \dot{r}.
  \end{aligned}
  \] Esta potencia la aporta la varilla: la fuerza de ligadura es
  perpendicular a la varilla, pero la varilla se mueve, y esa fuerza
  realiza trabajo sobre la cuenta.
\end{itemize}

\paragraph{\texorpdfstring{Oscilador con constante elástica variable:
\(H = E\), pero no se
conserva}{Oscilador con constante elástica variable: H = E, pero no se conserva}}\label{oscilador-con-constante-eluxe1stica-variable-h-e-pero-no-se-conserva}

Sea
\(L = \dfrac{1}{2} \cdot m \cdot \dot{x}^2 - \dfrac{1}{2} \cdot k(t) \cdot x^2\),
donde la constante elástica \(k(t)\) cambia con el tiempo por una acción
externa. La relación entre \(x\) y la posición no depende del tiempo,
así que \(T\) es homogénea de grado 2 y

\[
\begin{aligned}
H = \dfrac{p^2}{2 \cdot m} + \dfrac{1}{2} \cdot k(t) \cdot x^2 = T + V = E.
\end{aligned}
\]

Sin embargo,
\(\dfrac{dH}{dt} = \dfrac{\partial H}{\partial t} = \dfrac{1}{2} \cdot \dot{k}(t) \cdot x^2 \neq 0\):
la energía cambia porque el agente externo que modifica el muelle
realiza trabajo sobre el sistema.

\paragraph{Partícula cargada en un campo electromagnético: potencial que
depende de la
velocidad}\label{partuxedcula-cargada-en-un-campo-electromagnuxe9tico-potencial-que-depende-de-la-velocidad}

Para una partícula de carga \(q\) en un campo con potenciales
\(\phi(\mathbf{r},t)\) y \(\mathbf{A}(\mathbf{r},t)\), el Lagrangiano es

\[
\begin{aligned}
L = \dfrac{1}{2} \cdot m \cdot \mathbf{v} \cdot \mathbf{v} - q \cdot \phi + q \cdot \mathbf{v} \cdot \mathbf{A}.
\end{aligned}
\]

Aquí la hipótesis «\(V\) no depende de las velocidades» no se cumple:
hay un término lineal en las velocidades,
\(L_1 = q \cdot \mathbf{v} \cdot \mathbf{A}\). Aplicando
\(H = L_2 - L_0\):

\[
\begin{aligned}
H = \dfrac{1}{2} \cdot m \cdot \mathbf{v} \cdot \mathbf{v} + q \cdot \phi = \dfrac{\left|\mathbf{p} - q \cdot \mathbf{A}\right|^2}{2 \cdot m} + q \cdot \phi,
\end{aligned}
\]

donde \(\mathbf{p} = m \cdot \mathbf{v} + q \cdot \mathbf{A}\) es el
momento canónico. El término magnético no aparece en \(H\), lo que
concuerda con que la fuerza magnética no realiza trabajo. \(H\) es la
energía cinética más la energía potencial eléctrica, y se conserva solo
si \(\phi\) y \(\mathbf{A}\) no dependen explícitamente del tiempo.

\subsubsection{7.6. Resumen}\label{resumen}

Con un potencial independiente de las velocidades:

{\def\LTcaptype{none} % do not increment counter
\begin{longtable}[]{@{}
  >{\raggedright\arraybackslash}p{(\linewidth - 4\tabcolsep) * \real{0.3333}}
  >{\raggedright\arraybackslash}p{(\linewidth - 4\tabcolsep) * \real{0.3333}}
  >{\raggedright\arraybackslash}p{(\linewidth - 4\tabcolsep) * \real{0.3333}}@{}}
\toprule\noalign{}
\begin{minipage}[b]{\linewidth}\raggedright
Sistema
\end{minipage} & \begin{minipage}[b]{\linewidth}\raggedright
¿\(H = E\)?
\end{minipage} & \begin{minipage}[b]{\linewidth}\raggedright
¿Se conserva \(H\)?
\end{minipage} \\
\midrule\noalign{}
\endhead
\bottomrule\noalign{}
\endlastfoot
Esclerónomo, \(V(q)\) & Sí & Sí \\
Esclerónomo, \(V(q,t)\) & Sí & No \\
Reónomo, \(L\) sin \(t\) explícito & No: \(E - H = T_1 + 2 \cdot T_0\) &
Sí (integral de Jacobi) \\
Reónomo, \(L\) con \(t\) explícito & No & No \\
\end{longtable}
}

En mecánica cuántica, que toma la mecánica hamiltoniana como punto de
partida clásico, el operador Hamiltoniano \(\hat{H}\) se construye a
partir del Hamiltoniano clásico. En los sistemas del primer caso de la
tabla, que son los habituales en los problemas básicos (partícula en un
potencial, átomo de hidrógeno, oscilador armónico), \(\hat{H}\) es el
operador asociado a la energía total del sistema.

~

\end{document}
